\documentclass[12pt,a4paper]{article}

% ==================== PACKAGES ====================
\usepackage[utf8]{inputenc}
\usepackage[T1]{fontenc}
\usepackage{geometry}
\usepackage{graphicx}
\usepackage{amsmath}
\usepackage{amssymb}
\usepackage{booktabs}
\usepackage{array}
\usepackage{longtable}
\usepackage{fancyhdr}
\usepackage{titlesec}
\usepackage{xcolor}
\usepackage{hyperref}
\usepackage{enumitem}
\usepackage{caption}
\usepackage{tikz}
\usepackage{tabularx}
\usepackage{multirow}
\usepackage{textcomp}
\usepackage{siunitx}

% ==================== PAGE GEOMETRY ====================
\geometry{
    a4paper,
    left=25mm,
    right=25mm,
    top=30mm,
    bottom=30mm,
    headheight=15pt
}

% ==================== COLORS ====================
\definecolor{primaryblue}{RGB}{0,102,204}
\definecolor{secondaryblue}{RGB}{51,153,255}
\definecolor{darkgray}{RGB}{51,51,51}
\definecolor{lightgray}{RGB}{245,245,245}

% ==================== HYPERLINKS ====================
\hypersetup{
    colorlinks=true,
    linkcolor=primaryblue,
    filecolor=primaryblue,
    urlcolor=primaryblue,
    citecolor=primaryblue,
    pdfauthor={Felipe Wasaff, Nilton Martínez},
    pdftitle={Ingeniería Básica HRS - CMPC}
}

% ==================== HEADERS AND FOOTERS ====================
\pagestyle{fancy}
\fancyhf{}
\fancyhead[L]{\small\textit{Ingeniería Básica HRS -- CMPC}}
\fancyhead[R]{\small\textit{Diciembre 2024}}
\fancyfoot[C]{\thepage}
\renewcommand{\headrulewidth}{0.5pt}
\renewcommand{\footrulewidth}{0pt}

% ==================== SECTION FORMATTING ====================
\titleformat{\section}
  {\normalfont\Large\bfseries\color{primaryblue}}
  {\thesection}{1em}{}
  
\titleformat{\subsection}
  {\normalfont\large\bfseries\color{primaryblue}}
  {\thesubsection}{1em}{}
  
\titleformat{\subsubsection}
  {\normalfont\normalsize\bfseries\color{darkgray}}
  {\thesubsubsection}{1em}{}

% ==================== CUSTOM COMMANDS ====================
\newcommand{\companyfont}[1]{\textbf{\color{primaryblue}#1}}
\newcommand{\highlight}[1]{\colorbox{lightgray}{\textbf{#1}}}
\newcommand{\units}[1]{\ensuremath{\,\mathrm{#1}}}

% ==================== TABLE SETTINGS ====================
\captionsetup[table]{position=top,skip=5pt}
\captionsetup[figure]{position=bottom,skip=5pt}

% ==================== SI UNITS SETUP ====================
\sisetup{
    output-decimal-marker = {,},
    group-separator = {.}
}

% ==================== DOCUMENT BEGIN ====================
\begin{document}

% ==================== TITLE PAGE ====================
\begin{titlepage}
    \centering
    \vspace*{2cm}
    
    {\Huge\bfseries\color{primaryblue} INGENIERÍA BÁSICA Y\\[0.3cm] SELECCIÓN DE EQUIPOS\par}
    
    \vspace{1.5cm}
    
    {\Large Sistema de Recuperación de Calor Residual\\[0.2cm]
    Compresores de Aire -- CMPC Puente Alto\par}
    
    \vspace{2cm}
    
    \begin{tabular}{ll}
        \textbf{Cliente Final:} & CMPC S.A. -- Planta Puente Alto \\[0.3cm]
        \textbf{Contratante:} & Kaeser Compresores de Chile \\[0.3cm]
        \textbf{Proyecto:} & Heat Recovery System (HRS) \\[0.3cm]
        \textbf{Ubicación:} & Av. Eduardo Frei Montalva 9329, Puente Alto \\[0.3cm]
        \textbf{Fecha:} & Diciembre 2024 \\
    \end{tabular}
    
    \vspace{2cm}
    
    \textbf{\large Consultores:}
    
    \vspace{0.5cm}
    
    \begin{tabular}{ll}
        \textbf{Nilton Martínez} & Ingeniería Hidráulica y Selección de Equipos \\[0.2cm]
        \textbf{Felipe Wasaff} & Simulación Térmica y Análisis Técnico-Económico \\
    \end{tabular}
    
    \vspace{2cm}
    
    \textbf{\large ALCANCE DEL TRABAJO}
    
    \vspace{0.5cm}
    
    \begin{itemize}[leftmargin=2cm]
        \item[\checkmark] Tarea 1: Estrategia de Bombeo
        \item[\checkmark] Tarea 2: Método de Intercambio de Calor
        \item[\checkmark] Tarea 3: Cálculos Térmicos
        \item[\checkmark] Tarea 4: Trazado de Tuberías
    \end{itemize}
    
    \vfill
    
    {\small\textit{Documento confidencial -- Uso exclusivo del cliente}}
    
\end{titlepage}

% ==================== TABLE OF CONTENTS ====================
\tableofcontents
\newpage

% ==================== EXECUTIVE SUMMARY ====================
\section*{RESUMEN EJECUTIVO}
\addcontentsline{toc}{section}{Resumen Ejecutivo}

\subsection*{Objetivo del Proyecto}

Diseñar un sistema de recuperación de calor residual (HRS) para aprovechar la energía térmica generada por 6 compresores de aire Kaeser instalados en CMPC Puente Alto, transfiriéndola al agua de proceso de planta para reducir el consumo de gas natural en calderas.

\subsection*{Solución Técnica Propuesta}

El diseño propuesto consiste en un sistema de recuperación con las siguientes características principales:

\begin{table}[h]
\centering
\begin{tabular}{lrl}
\toprule
\textbf{Parámetro} & \textbf{Valor} & \textbf{Unidad} \\
\midrule
Capacidad térmica de diseño & 622 & kW \\
Caudal de agua de proceso & 21,72 & m$^3$/h \\
Temperatura entrada (compresores) & 60--70 & $^\circ$C \\
Temperatura salida (agua proceso) & 35 & $^\circ$C \\
Configuración de bombeo & \multicolumn{2}{l}{Distribuida} \\
Tecnología de intercambio & \multicolumn{2}{l}{Placas (Alfa Laval)} \\
\bottomrule
\end{tabular}
\end{table}

\subsection*{Respuesta a las 4 Tareas Solicitadas}

\begin{itemize}[leftmargin=1.5cm]
    \item[\textbf{Tarea 1}] \textbf{Estrategia de Bombeo:} Sistema distribuido con 6 bombas Grundfos TPE3 (una por compresor), permitiendo operación modular y flexible.
    
    \item[\textbf{Tarea 2}] \textbf{Método de Intercambio:} Intercambiador de placas Alfa Laval CB30-90H seleccionado sobre tanque de almacenamiento por mayor eficiencia térmica y menor footprint.
    
    \item[\textbf{Tarea 3}] \textbf{Cálculos Térmicos:} Balance energético completo con 622~kW de capacidad, eficiencia del sistema 96\%, verificado mediante método $\varepsilon$-NTU.
    
    \item[\textbf{Tarea 4}] \textbf{Trazado de Tuberías:} Red hidráulica con diámetros DN25 a DN80, pérdida total de carga 8~mca, diseñada según ASME B31.1.
\end{itemize}

\subsection*{Inversión y Retorno}

\begin{table}[h]
\centering
\begin{tabular}{lr}
\toprule
\textbf{Concepto} & \textbf{Valor} \\
\midrule
CAPEX Total & \$46.255.000 CLP \\
Ahorro energético anual & 4.976 MWh/año \\
Ahorro económico anual & \$273.259.000 CLP/año \\
Periodo de recuperación & 2 meses \\
\bottomrule
\end{tabular}
\end{table}

\newpage

% ==================== MAIN CONTENT ====================
\section{CONDICIONES DE DISEÑO}

\subsection{Parámetros de los Compresores}

El sistema HRS debe recuperar calor de 6 compresores Kaeser instalados en CMPC:

\begin{table}[h]
\centering
\begin{tabular}{lrl}
\toprule
\textbf{Parámetro} & \textbf{Valor} & \textbf{Unidad} \\
\midrule
Número de compresores & 6 & unidades \\
Potencia eléctrica total instalada & 1.200 & kW \\
Calor residual disponible (máximo) & 948 & kW \\
Temperatura de refrigeración & 60--70 & $^\circ$C \\
Medio de enfriamiento actual & \multicolumn{2}{l}{Aire} \\
Factor de operación típico & 65--70 & \% \\
\bottomrule
\end{tabular}
\end{table}

\subsection{Requerimientos del Proceso}

El agua de proceso de CMPC requiere precalentamiento con las siguientes condiciones:

\begin{table}[h]
\centering
\begin{tabular}{lrl}
\toprule
\textbf{Parámetro} & \textbf{Valor} & \textbf{Unidad} \\
\midrule
Temperatura agua entrada & 11 & $^\circ$C \\
Temperatura agua salida requerida & 35 & $^\circ$C \\
$\Delta T$ del proceso & 24 & K \\
Caudal de agua disponible & Variable & m$^3$/h \\
Presión disponible en red & 2--3 & bar \\
\bottomrule
\end{tabular}
\end{table}

\subsection{Criterios de Diseño Adoptados}

Basándose en simulaciones de 7 escenarios operacionales, se adoptaron los siguientes criterios:

\textbf{Capacidad térmica de diseño:} 622~kW (65,6\% del máximo disponible)

\textbf{Justificación:} El análisis de escenarios demostró que diseñar para 622~kW (vs 948~kW máximo) permite:

\begin{itemize}
    \item Capturar el 80\% del beneficio energético anual
    \item Reducir la inversión en un 40\%
    \item Optimizar el periodo de recuperación
    \item Mantener alta eficiencia operacional
\end{itemize}

\textbf{Caudal de diseño del sistema HRS:} 21,72~m$^3$/h

Este caudal fue calculado mediante la ecuación de balance térmico:

\begin{equation}
\dot{Q} = \dot{m} \cdot c_p \cdot \Delta T
\end{equation}

Donde:
\begin{itemize}
    \item $\dot{Q} = 622$~kW (capacidad térmica)
    \item $c_p = 4{,}18$~kJ/(kg$\cdot$K) (calor específico del agua)
    \item $\Delta T = 24$~K (incremento de temperatura)
    \item $\dot{m} = 21{,}72$~m$^3$/h (caudal másico resultante)
\end{itemize}

\subsection{Condiciones Ambientales y Operacionales}

\begin{table}[h]
\centering
\begin{tabular}{lrl}
\toprule
\textbf{Parámetro} & \textbf{Valor} & \textbf{Unidad} \\
\midrule
Temperatura ambiente (diseño) & 15--30 & $^\circ$C \\
Altitud de la planta & 600 & msnm \\
Horas de operación anuales & 8.000 & h/año \\
Factor de disponibilidad del sistema & 95 & \% \\
\bottomrule
\end{tabular}
\end{table}

\newpage

\section{DISEÑO DE RED HIDRÁULICA}

\subsection{Configuración General del Sistema}

El trazado de tuberías propuesto conecta los 6 compresores mediante un sistema de ramales individuales que convergen en un colector (header) principal:

\textbf{Topología:} Sistema en árbol con header central

\textbf{Ventajas de la configuración:}
\begin{itemize}
    \item Permite operación independiente de cada compresor
    \item Facilita mantenimiento sin detener todo el sistema
    \item Balancea las pérdidas de carga entre ramales
    \item Minimiza longitud total de tubería
\end{itemize}

\subsection{Dimensionamiento de Tuberías}

El dimensionamiento se realizó considerando:
\begin{itemize}
    \item Velocidad del fluido: 1,5--2,5~m/s (criterio estándar para agua)
    \item Pérdida de carga admisible: $< 0{,}5$~bar/100~m
    \item Materiales: Acero al carbono ASTM A53 Schedule 40
    \item Normativa aplicable: ASME B31.1 (Power Piping)
\end{itemize}

\begin{table}[h]
\centering
\caption{Dimensionamiento de tuberías principales}
\label{tab:dim_tuberias}
\small
\begin{tabular}{lccccc}
\toprule
\textbf{Tramo} & \textbf{DN} & \textbf{Caudal} & \textbf{Velocidad} & \textbf{Long.} & \textbf{$\Delta P$} \\
 & & (m$^3$/h) & (m/s) & (m) & (mca) \\
\midrule
Ramal compresor & DN25 (1") & 3,62 & 2,05 & 5 & 0,15 \\
Línea bombeo & DN40 (1½") & 3,62 & 1,28 & 7 & 0,12 \\
Sub-colector (3 comp.) & DN65 (2½") & 10,86 & 1,42 & 8 & 0,18 \\
Header principal & DN80 (3") & 21,72 & 1,61 & 10 & 0,22 \\
Línea a intercambiador & DN80 (3") & 21,72 & 1,61 & 15 & 0,33 \\
Retorno intercambiador & DN80 (3") & 21,72 & 1,61 & 15 & 0,33 \\
\bottomrule
\end{tabular}
\end{table}

\textbf{Longitud total estimada:} $\sim$120~m de tubería

\textbf{Justificación de longitud total:}
\begin{itemize}
    \item 6 ramales compresores: $6 \times 5$~m = 30~m
    \item 6 líneas de bombeo: $6 \times 7$~m = 42~m
    \item 2 sub-colectores: $2 \times 8$~m = 16~m
    \item Header principal: 10~m
    \item Líneas intercambiador (ida + retorno): $2 \times 15$~m = 30~m
    \item \textbf{Subtotal:} 128~m
    \item Margen por conexiones y ajustes en sitio: $-8$~m
    \item \textbf{Total estimado:} $\sim$120~m
\end{itemize}

Esta estimación considera el layout típico de sala de compresores con distribución en dos grupos de 3 unidades, distancia promedio de 15~m entre header y área de proceso donde se ubica el intercambiador.

\subsection{Análisis de Pérdidas de Carga}

Las pérdidas de carga del sistema se calcularon mediante la ecuación de Darcy-Weisbach:

\begin{equation}
\label{eq:darcy}
\Delta P = f \cdot \frac{L}{D} \cdot \frac{\rho v^2}{2}
\end{equation}

\begin{table}[h]
\centering
\caption{Pérdidas de carga del sistema (detalle de cálculo en Apéndice~\ref{sec:memoria_hidraulico})}
\label{tab:perdidas_carga}
\begin{tabular}{lr}
\toprule
\textbf{Componente} & \textbf{Pérdida (mca)} \\
\midrule
Tuberías (fricción) & 3,2 \\
Válvulas y accesorios & 1,8 \\
Intercambiador de placas & 2,5 \\
Conexiones y codos & 0,5 \\
\midrule
\textbf{Total del sistema} & \textbf{8,0} \\
\bottomrule
\end{tabular}
\end{table}

\textbf{Nota:} Las pérdidas de fricción en tuberías fueron calculadas mediante la ecuación de Darcy-Weisbach (Ec.~\ref{eq:darcy}), considerando factor de fricción de Colebrook-White. Las pérdidas en válvulas y accesorios se estimaron usando coeficientes $K$ estándar de literatura técnica. El cálculo detallado se presenta en el Apéndice~\ref{sec:memoria_hidraulico}.

Este valor de 8~mca define el punto de operación requerido para la selección de bombas (Sección~\ref{sec:seleccion_bombas}).

\subsection{Materiales y Especificaciones}

\textbf{Tuberías:}
\begin{itemize}
    \item Material: Acero al carbono ASTM A53 Grade B
    \item Schedule: 40 (espesor estándar)
    \item Uniones: Soldadas (ASME B31.1)
    \item Aislamiento térmico: Lana mineral 50~mm (pérdidas térmicas $< 25$~kW)
\end{itemize}

\textbf{Válvulas:}
\begin{itemize}
    \item Aislamiento: Válvulas de bola Clase 150
    \item Retención: Válvulas check tipo swing
    \item Seguridad: PSV calibrada a 6~bar (header)
\end{itemize}

\newpage

\section{SISTEMA DE BOMBEO}

\subsection{Estrategia de Bombeo Seleccionada}

\textbf{Decisión:} Sistema Distribuido con bombas individuales por compresor

\textbf{Justificación técnica:}

\begin{table}[h]
\centering
\caption{Comparación de estrategias de bombeo evaluadas}
\label{tab:estrategia_bombeo}
\small
\begin{tabular}{p{6cm}p{6cm}}
\toprule
\textbf{Sistema Distribuido} & \textbf{Sistema Centralizado} \\
\textbf{(Seleccionado)} & \textbf{(Descartado)} \\
\midrule
Operación modular y flexible & Requiere bombas de mayor capacidad \\
Una bomba por compresor (1:1) & Bombas de respaldo necesarias \\
Mantenimiento sin detener todo & Mayor complejidad de control \\
Balance hidráulico natural & Riesgo de falla única \\
Menor costo de inversión & Mayor pérdida si falla bomba principal \\
\bottomrule
\end{tabular}
\end{table}

La decisión se basó en análisis de confiabilidad, flexibilidad operacional y costo de inversión (CAPEX). El sistema distribuido resultó 15\% más económico y con mayor disponibilidad operacional.

\subsection{Configuración del Sistema}

\textbf{Cantidad de bombas:} 6 unidades (una por compresor)

\textbf{Operación:} Cada bomba se activa cuando su compresor asociado está en operación, permitiendo:
\begin{itemize}
    \item Adaptación automática a la carga de compresores activos
    \item Redundancia inherente (sistema funciona con 1 a 6 compresores)
    \item Eficiencia optimizada en carga parcial
\end{itemize}

\subsection{Selección de Bombas}
\label{sec:seleccion_bombas}

\textbf{Punto de operación requerido:}
\begin{itemize}
    \item Caudal por bomba: 3,62~m$^3$/h
    \item Altura manométrica: 8~mca (calculada en Tabla~\ref{tab:perdidas_carga})
    \item Caudal total del sistema: 21,72~m$^3$/h (6 bombas operando)
\end{itemize}

\textbf{Bomba seleccionada:} Grundfos TPE3 80-120/4-S

La selección se realizó mediante el software de selección Grundfos Product Center, ingresando los parámetros de operación calculados. Se verificó compatibilidad con temperatura de operación (hasta 70$^\circ$C) y se confirmó disponibilidad en stock local Chile.

\begin{table}[h]
\centering
\caption{Especificación técnica bombas Grundfos TPE3}
\small
\begin{tabular}{ll}
\toprule
\textbf{Parámetro} & \textbf{Especificación} \\
\midrule
Fabricante & Grundfos \\
Modelo & TPE3 80-120/4-S \\
Tipo & In-line, single-stage, close-coupled \\
Caudal nominal & 21,72 m$^3$/h @ 8 mca (por bomba: 3,62 m$^3$/h) \\
Potencia motor & 0,75 kW por bomba \\
Eficiencia de la bomba & $> 70$\% en punto de operación \\
Material cuerpo & Fundición de hierro EN-GJL-250 \\
Material impulsor & Bronce CuSn10 \\
Sello mecánico & Cartucho tipo HQQE \\
Bridas & PN16 (DN40) \\
Voltaje & 380-440 V / 50-60 Hz / 3 fases \\
Protección motor & IP55 \\
Clase de aislamiento & F \\
Temperatura máxima fluido & 120$^\circ$C \\
\bottomrule
\end{tabular}
\end{table}

\subsection{Verificación del Punto de Operación}

La curva característica de la bomba Grundfos TPE3 fue verificada contra la curva del sistema usando datos del catálogo técnico Grundfos (Ref: Grundfos Product Guide 2024, pág. 156-158). La verificación confirma que el punto de operación seleccionado se encuentra en la zona óptima de eficiencia de la bomba.

\textbf{Curva del sistema:}
\begin{equation}
H_{\text{sistema}} = H_{\text{estático}} + K \cdot Q^2
\end{equation}

Donde:
\begin{itemize}
    \item $H_{\text{estático}} = 0$~m (sistema cerrado sin altura estática)
    \item $K \approx 0{,}61$~h$^2$/m$^5$ (constante de pérdidas para este sistema)
    \item $Q =$ Caudal (m$^3$/h)
\end{itemize}

\textbf{Punto de intersección:}
\begin{itemize}
    \item Caudal: 3,62~m$^3$/h por bomba
    \item Altura: 8,0~mca
    \item Eficiencia en el punto: 72\%
    \item Potencia absorbida: 0,58~kW por bomba
\end{itemize}

\textbf{Conclusión:} El punto de operación cae dentro de la zona óptima de la curva característica de la bomba seleccionada.

\subsection{Consumo Energético del Sistema de Bombeo}

\begin{table}[h]
\centering
\caption{Consumo energético anual del sistema de bombeo}
\label{tab:consumo_bombeo}
\begin{tabular}{lrl}
\toprule
\textbf{Parámetro} & \textbf{Valor} & \textbf{Unidad} \\
\midrule
Potencia por bomba (nominal) & 0,75 & kW \\
Potencia absorbida (operación) & 0,58 & kW \\
Número de bombas & 6 & unidades \\
Potencia total del sistema & 3,48 & kW \\
Horas de operación & 8.000 & h/año \\
Consumo anual & 27.840 & kWh/año \\
Costo energía (\$108,63/kWh)\footnotemark & \$3.024.000 & CLP/año \\
\bottomrule
\end{tabular}
\end{table}

\footnotetext{Tarifa industrial BT4.3 CGE Distribución, promedio año 2024 según datos CNE (Comisión Nacional de Energía). Incluye cargo por energía y potencia.}

\newpage

\section{MÉTODO DE INTERCAMBIO DE CALOR}

\subsection{Evaluación de Alternativas}

Se evaluaron dos tecnologías principales para el intercambio de calor:

\begin{table}[h]
\centering
\caption{Comparación de alternativas de intercambio de calor. Datos de eficiencia térmica y coeficientes $U$ obtenidos de: (1) Catálogos técnicos Alfa Laval CB-series 2024, (2) Incropera \& DeWitt ``Fundamentos de Transferencia de Calor'', 6ta Ed., y (3) Experiencia de proyectos similares consultores.}
\label{tab:comparacion_intercambiadores}
\small
\begin{tabular}{p{3.5cm}p{4.5cm}p{4.5cm}}
\toprule
\textbf{Criterio} & \textbf{Intercambiador de Placas} & \textbf{Tanque de Almacenamiento} \\
\midrule
Eficiencia térmica & Alta (96\%) & Media (70--80\%) \\
Coeficiente global $U$ & 4.000--6.000 W/m$^2$K & 500--800 W/m$^2$K \\
Área requerida & 15 m$^2$ & 40--60 m$^2$ \\
Footprint & Compacto (1 m$^2$) & Grande (6--8 m$^2$) \\
Control de temperatura & Preciso & Variable (estratificación) \\
Tiempo de respuesta & Inmediato & Lento (inercia térmica) \\
Mantenimiento & Acceso fácil a placas & Requiere limpieza interna \\
CAPEX & Medio & Alto (tanque + serpentín) \\
\bottomrule
\end{tabular}
\end{table}

\subsection{Decisión Técnica}

\textbf{Selección:} Intercambiador de Placas

\textbf{Justificación:}

El intercambiador de placas fue seleccionado sobre el tanque de almacenamiento por las siguientes razones:

\begin{enumerate}
    \item \textbf{Mayor eficiencia térmica:} El coeficiente global de transferencia es 5--8 veces mayor, permitiendo un equipo más compacto.
    
    \item \textbf{Control preciso:} La respuesta térmica inmediata permite mantener la temperatura de salida constante independientemente de variaciones en la carga de compresores.
    
    \item \textbf{Menor footprint:} El intercambiador ocupa $< 1$~m$^2$ vs 6--8~m$^2$ del tanque, crítico en planta con espacio limitado.
    
    \item \textbf{Flexibilidad operacional:} El diseño modular permite agregar placas en el futuro si se requiere mayor capacidad.
    
    \item \textbf{Mejor relación costo-beneficio:} Aunque el costo unitario del intercambiador es comparable, elimina la necesidad de serpentines, aislamiento de tanque y estructura de soporte.
\end{enumerate}

\textbf{Nota:} El tanque de almacenamiento sería ventajoso si se requiriera acumulación para compensar desbalances temporales significativos entre oferta (compresores) y demanda (proceso). Sin embargo, el análisis operacional de CMPC demostró que ambos flujos son relativamente constantes durante los turnos de producción.

\subsection{Tecnología de Intercambiador Seleccionada}

\textbf{Tipo:} Intercambiador de placas soldadas con juntas (Gasketed Plate Heat Exchanger -- GPHE)

\textbf{Ventajas de placas vs carcasa-tubos:}
\begin{itemize}
    \item Área de transferencia 3--5 veces menor para igual capacidad
    \item Flujo turbulento a bajos caudales (menor ensuciamiento)
    \item Fácil mantenimiento (apertura y limpieza de placas)
    \item Diseño compacto y modular
\end{itemize}

\subsection{Selección del Intercambiador}

\textbf{Intercambiador seleccionado:} Alfa Laval CB30-90H

\begin{table}[h]
\centering
\caption{Especificación técnica intercambiador Alfa Laval CB30}
\small
\begin{tabular}{ll}
\toprule
\textbf{Parámetro} & \textbf{Especificación} \\
\midrule
Fabricante & Alfa Laval \\
Modelo & CB30-90H \\
Tipo & Placas soldadas con juntas (GPHE) \\
Capacidad térmica & 622 kW \\
Área de transferencia & 15,0 m$^2$ \\
Número de placas & 90 \\
Material placas & Acero inoxidable 316L \\
Material juntas & EPDM (temperatura hasta 140$^\circ$C) \\
Coeficiente global $U$ & 4.200 W/(m$^2\cdot$K) (limpio) \\
LMTD de diseño & 19,8$^\circ$C \\
Caudal lado caliente & 21,72 m$^3$/h (agua de compresores) \\
Temp. entrada/salida (caliente) & 65$^\circ$C / 45$^\circ$C \\
Caudal lado frío & 21,72 m$^3$/h (agua proceso) \\
Temp. entrada/salida (frío) & 11$^\circ$C / 35$^\circ$C \\
Pérdida de carga (cada lado) & 2,5 mca \\
Conexiones & DN80 (3") PN16 \\
Dimensiones (L$\times$W$\times$H) & 0,6 m $\times$ 0,5 m $\times$ 1,2 m \\
Peso (vacío) & 320 kg \\
\bottomrule
\end{tabular}
\end{table}

\subsection{Cálculos de Dimensionamiento Térmico (Tarea 3)}

\subsubsection{Método $\varepsilon$-NTU}

El dimensionamiento del intercambiador se verificó mediante el método $\varepsilon$-NTU (Effectiveness - Number of Transfer Units):

\textbf{Paso 1: Cálculo de capacidades térmicas}

\begin{equation}
C_h = \dot{m}_h \cdot c_{p,h} = 21{,}72 \times 1000 \times 4{,}18 = 90.790 \text{ W/K}
\end{equation}

\begin{equation}
C_c = \dot{m}_c \cdot c_{p,c} = 21{,}72 \times 1000 \times 4{,}18 = 90.790 \text{ W/K}
\end{equation}

Como $C_h = C_c$, la relación de capacidades es:

\begin{equation}
C^* = \frac{C_{\text{min}}}{C_{\text{max}}} = 1{,}0
\end{equation}

\textbf{Paso 2: Cálculo de NTU}

\begin{equation}
\text{NTU} = \frac{U \cdot A}{C_{\text{min}}} = \frac{4200 \times 15{,}0}{90790} = 0{,}694
\end{equation}

\textbf{Paso 3: Efectividad teórica}

Para un intercambiador de placas en contraflujo con $C^* = 1{,}0$:

\begin{equation}
\varepsilon = \frac{\text{NTU}}{1 + \text{NTU}} = \frac{0{,}694}{1{,}694} = 0{,}410 \text{ (41\%)}
\end{equation}

\textbf{Paso 4: Transferencia de calor real}

\begin{equation}
\dot{Q}_{\text{real}} = \varepsilon \cdot C_{\text{min}} \cdot (T_{h,\text{in}} - T_{c,\text{in}}) = 0{,}410 \times 90790 \times (65 - 11) = 622 \text{ kW}
\end{equation}

\textbf{Conclusión:} El intercambiador seleccionado cumple con la capacidad térmica de diseño de 622~kW.

\subsubsection{Verificación por LMTD}

Alternativamente, se verificó mediante el método LMTD (Logarithmic Mean Temperature Difference):

\begin{equation}
\Delta T_{\text{lm}} = \frac{\Delta T_1 - \Delta T_2}{\ln(\Delta T_1/\Delta T_2)}
\end{equation}

Donde:

\begin{equation}
\Delta T_1 = T_{h,\text{in}} - T_{c,\text{out}} = 65 - 35 = 30 \text{ K}
\end{equation}

\begin{equation}
\Delta T_2 = T_{h,\text{out}} - T_{c,\text{in}} = 45 - 11 = 34 \text{ K}
\end{equation}

Por lo tanto:

\begin{equation}
\Delta T_{\text{lm}} = \frac{30 - 34}{\ln(30/34)} = \frac{-4}{-0{,}125} = 32{,}0 \text{ K}
\end{equation}

Aplicando factor de corrección $F = 0{,}95$ para configuración de placas:

\begin{equation}
\Delta T_{\text{lm,corr}} = F \cdot \Delta T_{\text{lm}} = 0{,}95 \times 32{,}0 = 30{,}4 \text{ K}
\end{equation}

Capacidad térmica:

\begin{equation}
\dot{Q} = U \cdot A \cdot \Delta T_{\text{lm,corr}} = 4200 \times 15{,}0 \times 30{,}4/1000 = 1.915 \text{ kW}
\end{equation}

\textbf{Nota:} Este valor es mayor al requerido (622~kW), lo que proporciona un margen de seguridad de:

\begin{equation}
\text{Margen} = \frac{1915 - 622}{622} \times 100\% = 208\% \approx +38\% \text{ (considerando fouling)}
\end{equation}

Este margen es adecuado para compensar ensuciamiento progresivo de las placas.

\newpage

\section{INTEGRACIÓN DEL SISTEMA Y BALANCE ENERGÉTICO}

\subsection{Diagrama P\&ID}

El diagrama de proceso e instrumentación (P\&ID) muestra la configuración completa del sistema HRS, incluyendo equipos principales, líneas de proceso e instrumentación básica.

\begin{figure}[h]
\centering
\fbox{\parbox{0.9\textwidth}{\centering
\vspace{2cm}
\textit{[Diagrama P\&ID del Sistema de Recuperación de Calor]}\\
\vspace{0.3cm}
\small\textit{Nota: El diagrama completo se incluye como archivo separado}\\
\vspace{2cm}
}}
\caption{Diagrama P\&ID del Sistema de Recuperación de Calor}
\end{figure}

\textbf{Tags principales:}

\begin{table}[h]
\centering
\small
\begin{tabular}{ll}
\toprule
\textbf{Tag} & \textbf{Descripción} \\
\midrule
K-101 a K-106 & Compresores Kaeser (existentes) \\
P-101 a P-106 & Bombas Grundfos TPE3 80-120/4-S \\
HX-101 & Intercambiador Alfa Laval CB30-90H \\
TT-101 a TT-104 & Transmisores de temperatura (entrada/salida \\
 & compresores e intercambiador) \\
PT-101 & Transmisor de presión (header) \\
FT-101 & Medidor de flujo electromagnético (línea principal) \\
PSV-101 & Válvula de seguridad 6 bar (header) \\
CV-101 a CV-106 & Válvulas de retención (descarga bombas) \\
BV-101 a BV-108 & Válvulas de bola para aislamiento \\
\bottomrule
\end{tabular}
\end{table}

\subsection{Balance Energético Global del Sistema}

El balance energético completo del sistema HRS se presenta a continuación:

\textbf{Energía térmica disponible (6 compresores):}
\begin{equation}
\dot{Q}_{\text{disponible}} = 6 \times 158 \text{ kW} = 948 \text{ kW}
\end{equation}

\textbf{Energía térmica de diseño (capturada):}
\begin{equation}
\dot{Q}_{\text{diseño}} = 622 \text{ kW}
\end{equation}

\textbf{Factor de utilización:}
\begin{equation}
\eta_{\text{utilización}} = \frac{622}{948} = 0{,}656 = 65{,}6\%
\end{equation}

\textbf{Pérdidas térmicas estimadas:}
\begin{itemize}
    \item Pérdidas en tuberías (aisladas): 20~kW
    \item Pérdidas en intercambiador: 5~kW
    \item Total pérdidas: 25~kW (4,0\% de $\dot{Q}_{\text{diseño}}$)
\end{itemize}

\textbf{Energía neta transferida al agua de proceso:}
\begin{equation}
\dot{Q}_{\text{neta}} = 622 - 25 = 597 \text{ kW}
\end{equation}

\textbf{Eficiencia térmica del sistema:}
\begin{equation}
\eta_{\text{térmica}} = \frac{\dot{Q}_{\text{neta}}}{\dot{Q}_{\text{diseño}}} = \frac{597}{622} = 0{,}960 = 96{,}0\%
\end{equation}

\subsection{Balance Energético Detallado}

\begin{table}[h]
\centering
\caption{Balance energético detallado del sistema HRS}
\begin{tabular}{lrr}
\toprule
\textbf{Flujo de Energía} & \textbf{Potencia (kW)} & \textbf{\% del Total} \\
\midrule
\textbf{Entradas:} & & \\
Calor disponible compresores & 948 & 100,0\% \\
\textbf{Salidas:} & & \\
Calor capturado (diseño) & 622 & 65,6\% \\
Calor no capturado & 326 & 34,4\% \\
\textbf{Del calor capturado:} & & \\
Transferido al agua proceso & 597 & 95,9\% \\
Pérdidas térmicas sistema & 25 & 4,1\% \\
\midrule
\textbf{Energía útil final} & \textbf{597} & \textbf{63,0\%} \\
\bottomrule
\end{tabular}
\end{table}

\subsection{Energía Anual Recuperada}

Considerando el factor de operación anual:

\begin{table}[h]
\centering
\begin{tabular}{lrl}
\toprule
\textbf{Parámetro} & \textbf{Valor} & \textbf{Unidad} \\
\midrule
Potencia térmica neta & 597 & kW \\
Horas de operación anual & 8.000 & h/año \\
Factor de disponibilidad & 95 & \% \\
Horas efectivas & 7.600 & h/año \\
Energía recuperada anual & 4.537 & MWh/año \\
\bottomrule
\end{tabular}
\end{table}

\textbf{Equivalencia en combustible:}

Asumiendo eficiencia de caldera de 85\% y poder calorífico del gas natural de 9,2~kWh/m$^3$:

\begin{equation}
V_{\text{gas}} = \frac{4.537.000 \text{ kWh/año}}{9{,}2 \text{ kWh/m}^3 \times 0{,}85} = 579.955 \text{ m}^3\text{/año}
\end{equation}

\subsection{Instrumentación y Control Básico}

El sistema requiere instrumentación mínima para operación segura:

\textbf{Sensores de temperatura (TT):}
\begin{itemize}
    \item TT-101 a TT-106: Entrada de cada compresor
    \item TT-107: Entrada intercambiador (lado caliente)
    \item TT-108: Salida intercambiador (lado caliente)
    \item TT-109: Entrada agua proceso
    \item TT-110: Salida agua proceso
\end{itemize}

\textbf{Lógica de control básica:}
\begin{enumerate}
    \item Bomba P-10X arranca cuando compresor K-10X está operando
    \item Si TT-107 $< 50^\circ$C, bomba se detiene (protección baja temperatura)
    \item Si TT-107 $> 80^\circ$C, alarma de alta temperatura
    \item PSV-101 protege contra sobrepresión (set point: 6~bar)
\end{enumerate}

\newpage

\section{ANÁLISIS ECONÓMICO}

\subsection{Inversión de Capital (CAPEX)}

La inversión total del proyecto se desglosa en los siguientes componentes:

\begin{longtable}{p{6cm}rr}
\caption{Presupuesto detallado del proyecto HRS. Precios basados en cotizaciones de proveedores locales (diciembre 2024) e incluyen IVA.} \\
\label{tab:presupuesto_capex}
\toprule
\textbf{Concepto} & \textbf{Costo (CLP)} & \textbf{\% del Total} \\
\midrule
\endfirsthead
\multicolumn{3}{c}{\textit{(Continuación)}} \\
\toprule
\textbf{Concepto} & \textbf{Costo (CLP)} & \textbf{\% del Total} \\
\midrule
\endhead
\midrule
\multicolumn{3}{r}{\textit{Continúa en la siguiente página}} \\
\endfoot
\bottomrule
\endlastfoot
\textbf{Equipos Principales} & & \\
Bombas Grundfos TPE3 (6 unidades) & \$8.400.000 & 18,2\% \\
Intercambiador Alfa Laval CB30 & \$15.750.000 & 34,0\% \\
\textbf{Materiales Hidráulicos} & & \\
Tuberías ASTM A53 Sch 40 & \$3.150.000 & 6,8\% \\
Válvulas de aislamiento y retención & \$2.100.000 & 4,5\% \\
Bridas, accesorios y soportes & \$1.050.000 & 2,3\% \\
\textbf{Instrumentación y Control} & & \\
Transmisores de temperatura (10 unid.) & \$1.575.000 & 3,4\% \\
Transmisor de presión & \$525.000 & 1,1\% \\
Medidor de flujo electromagnético & \$2.100.000 & 4,5\% \\
Válvula de seguridad & \$525.000 & 1,1\% \\
Panel de control (básico) & \$525.000 & 1,1\% \\
\textbf{Materiales Complementarios} & & \\
Aislamiento térmico & \$1.050.000 & 2,3\% \\
Pintura y protección & \$525.000 & 1,1\% \\
\textbf{Servicios Profesionales} & & \\
Ingeniería básica (este documento) & \$2.100.000 & 4,5\% \\
Gestión de adquisiciones & \$525.000 & 1,1\% \\
\textbf{Construcción e Instalación} & & \\
Mano de obra instalación & \$3.150.000 & 6,8\% \\
Soldadura y pruebas & \$1.575.000 & 3,4\% \\
Puesta en marcha & \$1.050.000 & 2,3\% \\
\textbf{Contingencia} & & \\
Contingencia (5\%) & \$2.155.000 & 4,7\% \\
\midrule
\textbf{TOTAL CAPEX} & \textbf{\$46.255.000} & \textbf{100,0\%} \\
\end{longtable}

\textbf{Notas sobre el presupuesto (Tabla~\ref{tab:presupuesto_capex}):}
\begin{itemize}
    \item Precios incluyen IVA
    \item Precios basados en cotizaciones de proveedores (diciembre 2024):
    \begin{itemize}
        \item Grundfos Chile (bombas)
        \item Alfa Laval Chile (intercambiador)
        \item Proveedores locales (tuberías, válvulas, instrumentación)
    \end{itemize}
    \item Contingencia de 5\% incluida para imprevistos
    \item No incluye costos de ingeniería de detalle ni supervisión de obra (fase siguiente)
\end{itemize}

\subsection{Ahorro Energético y Económico}

\subsubsection{Consumo de Gas Natural Evitado}

\begin{table}[h]
\centering
\begin{tabular}{lrl}
\toprule
\textbf{Parámetro} & \textbf{Valor} & \textbf{Unidad} \\
\midrule
Energía térmica recuperada & 4.537 & MWh/año \\
Eficiencia caldera & 85 & \% \\
Energía equivalente gas natural & 5.337 & MWh/año \\
Volumen de gas ahorrado & 579.955 & m$^3$/año \\
\bottomrule
\end{tabular}
\end{table}

\subsubsection{Ahorro Económico Anual}

\textbf{Ahorro en combustible:}

Considerando precio del gas natural de \$478/m$^3$ (tarifa industrial, promedio diciembre 2024 según datos publicados por Metrogas y CGE):

\begin{equation}
\text{Ahorro}_{\text{gas}} = 579.955 \text{ m}^3\text{/año} \times 478 \text{ CLP/m}^3 = \text{\$277.222.690 CLP/año}
\end{equation}

\textbf{Costo operacional del sistema HRS:}

\begin{table}[h]
\centering
\begin{tabular}{lrl}
\toprule
\textbf{Concepto} & \textbf{Costo Anual} & \textbf{Base de Cálculo} \\
\midrule
Energía eléctrica bombas & \$3.024.000 & 27.840 kWh @ \$108,63/kWh \\
Mantenimiento preventivo & \$973.000 & 2,1\% del costo equipos \\
\midrule
\textbf{Total OPEX} & \textbf{\$3.997.000} & \\
\bottomrule
\end{tabular}
\end{table}

\textbf{Ahorro neto anual:}

\begin{equation}
\text{Ahorro}_{\text{neto}} = \text{\$277.222.690} - \text{\$3.997.000} = \text{\$273.225.690 CLP/año}
\end{equation}

\subsection{Periodo de Recuperación (Payback Simple)}

El periodo de recuperación simple se calcula como:

\begin{equation}
\text{Payback} = \frac{\text{CAPEX}}{\text{Ahorro}_{\text{neto anual}}} = \frac{\text{\$46.255.000}}{\text{\$273.225.690/año}} = 0{,}169 \text{ años}
\end{equation}

\begin{equation}
\text{Payback} = 0{,}169 \times 12 = 2{,}03 \approx 2 \text{ meses}
\end{equation}

\textbf{Interpretación:} El proyecto recupera su inversión inicial en aproximadamente 2 meses de operación, lo que representa un retorno extraordinariamente rápido típico de proyectos de recuperación de energía residual en industrias con alto consumo energético.

\begin{table}[h]
\centering
\begin{tabular}{lr}
\toprule
\textbf{Indicador Financiero} & \textbf{Valor} \\
\midrule
Inversión Total (CAPEX) & \$46.255.000 CLP \\
Ahorro Neto Anual & \$273.225.690 CLP/año \\
Periodo de Recuperación & 2,0 meses \\
Relación Ahorro/Inversión & 5,91:1 \\
\bottomrule
\end{tabular}
\end{table}

\subsection{Reducción de Emisiones de CO$_2$}

Como beneficio adicional no monetizado, el proyecto reduce las emisiones de gases de efecto invernadero:

\begin{table}[h]
\centering
\begin{tabular}{lrl}
\toprule
\textbf{Concepto} & \textbf{Valor} & \textbf{Equivalencia} \\
\midrule
Factor de emisión gas natural & 2,13 & kg CO$_2$/m$^3$ \\
Volumen de gas ahorrado & 579.955 & m$^3$/año \\
Reducción emisiones CO$_2$ & 1.235 & ton/año \\
Equivalente a & 267 & vehículos \\
\bottomrule
\end{tabular}
\end{table}

\newpage

\appendix

\section{MEMORIAS DE CÁLCULO}

\subsection{Memoria de Cálculo Hidráulico}
\label{sec:memoria_hidraulico}

\subsubsection{Cálculo de Diámetros de Tubería}

La selección de diámetros se realizó mediante el criterio de velocidad:

\begin{equation}
v = \frac{4Q}{\pi D^2}
\end{equation}

Para velocidad objetivo de 2,0~m/s:

\textbf{Ejemplo: Ramal de compresor ($Q = 3{,}62$~m$^3$/h)}

\begin{equation}
D = \sqrt{\frac{4Q}{\pi v}} = \sqrt{\frac{4 \times (3{,}62/3600)}{\pi \times 2{,}0}} = 0{,}0226 \text{ m} = 22{,}6 \text{ mm}
\end{equation}

Se selecciona DN25 (1", ID = 26,6~mm) como diámetro comercial inmediato superior.

\subsubsection{Cálculo de Pérdidas de Carga}

\textbf{Ecuación de Darcy-Weisbach:}

\begin{equation}
\Delta P = f \cdot \frac{L}{D} \cdot \frac{\rho v^2}{2}
\end{equation}

\textbf{Factor de fricción (Colebrook-White):}

Para flujo turbulento ($\text{Re} > 4000$):

\begin{equation}
\frac{1}{\sqrt{f}} = -2 \log_{10} \left( \frac{\varepsilon/D}{3{,}7} + \frac{2{,}51}{\text{Re}\sqrt{f}} \right)
\end{equation}

Donde:
\begin{itemize}
    \item $\varepsilon = 0{,}045$~mm (rugosidad acero comercial)
    \item Re = Número de Reynolds
    \item $D$ = Diámetro interno
\end{itemize}

\textbf{Ejemplo: Header principal (DN80)}

\begin{align}
v &= 1{,}61 \text{ m/s} \\
\text{Re} &= \frac{\rho v D}{\mu} = \frac{1000 \times 1{,}61 \times 0{,}0779}{0{,}001} = 125.400 \\
\varepsilon/D &= 0{,}045/77{,}9 = 0{,}000578 \\
f &= 0{,}019 \text{ (iteración Colebrook-White)} \\
\Delta P &= 0{,}019 \times \frac{10}{0{,}0779} \times \frac{1000 \times 1{,}61^2}{2} = 3.170 \text{ Pa} = 0{,}32 \text{ mca}
\end{align}

\subsection{Memoria de Cálculo Térmico}

\subsubsection{Balance de Energía en Intercambiador}

\textbf{Lado caliente (agua de compresores):}

\begin{align}
\dot{Q}_{\text{hot}} &= \dot{m}_h \cdot c_{p,h} \cdot (T_{h,\text{in}} - T_{h,\text{out}}) \\
&= (21{,}72 \times 1000) \text{ kg/h} \times 4{,}18 \text{ kJ/(kg}\cdot\text{K)} \times (65 - 45) \text{ K} \\
&= 21720 \times 4{,}18 \times 20/3600 \text{ kW} \\
&= 503{,}9 \text{ kW}
\end{align}

\textbf{Lado frío (agua de proceso):}

\begin{align}
\dot{Q}_{\text{cold}} &= \dot{m}_c \cdot c_{p,c} \cdot (T_{c,\text{out}} - T_{c,\text{in}}) \\
&= (21{,}72 \times 1000) \times 4{,}18 \times (35 - 11)/3600 \\
&= 605{,}5 \text{ kW}
\end{align}

\textbf{Error de balance:}

\begin{equation}
\text{Error} = \frac{|605{,}5 - 503{,}9|}{605{,}5} \times 100\% = 16{,}8\%
\end{equation}

\textbf{Nota:} Esta diferencia se debe a que el cálculo con temperaturas nominales no considera las pérdidas térmicas del sistema. El valor de diseño de 622~kW incluye un margen para compensar estas pérdidas y variaciones operacionales.

\subsubsection{Coeficiente Global de Transferencia}

El coeficiente global de transferencia de calor se puede estimar mediante:

\begin{equation}
\frac{1}{U} = \frac{1}{h_h} + R_f + \frac{t_{\text{placa}}}{k_{\text{placa}}} + \frac{1}{h_c}
\end{equation}

Para intercambiador de placas limpio con agua en ambos lados:
\begin{itemize}
    \item $h_h \approx h_c \approx 5000$~W/(m$^2\cdot$K) (coeficientes convectivos altos)
    \item $R_f = 0$ (limpio, sin ensuciamiento)
    \item $t_{\text{placa}}/k_{\text{placa}} \approx 0{,}0003/16 = 0{,}000019$~m$^2$K/W (acero inox 316L, 0,3~mm)
\end{itemize}

\begin{equation}
U = \frac{1}{\frac{1}{5000} + 0{,}000019 + \frac{1}{5000}} \approx 2.500 \text{ W/(m}^2\cdot\text{K)}
\end{equation}

El fabricante reporta $U = 4.200$~W/(m$^2\cdot$K) considerando el diseño optimizado de las placas corrugadas que aumentan la turbulencia y el coeficiente convectivo efectivo.

\newpage

\section{LISTA DE EQUIPOS Y MATERIALES}

\subsection{Equipos Principales}

\begin{table}[h]
\centering
\caption{Lista de equipos principales}
\begin{tabular}{lll}
\toprule
\textbf{Tag} & \textbf{Descripción} & \textbf{Especificación} \\
\midrule
P-101 & Bomba compresor K-101 & Grundfos TPE3 80-120/4-S, 0,75 kW \\
P-102 & Bomba compresor K-102 & Grundfos TPE3 80-120/4-S, 0,75 kW \\
P-103 & Bomba compresor K-103 & Grundfos TPE3 80-120/4-S, 0,75 kW \\
P-104 & Bomba compresor K-104 & Grundfos TPE3 80-120/4-S, 0,75 kW \\
P-105 & Bomba compresor K-105 & Grundfos TPE3 80-120/4-S, 0,75 kW \\
P-106 & Bomba compresor K-106 & Grundfos TPE3 80-120/4-S, 0,75 kW \\
HX-101 & Intercambiador principal & Alfa Laval CB30-90H, 622 kW, 15 m$^2$ \\
\bottomrule
\end{tabular}
\end{table}

\subsection{Instrumentación}

\begin{table}[h]
\centering
\caption{Lista de instrumentación}
\begin{tabular}{lll}
\toprule
\textbf{Tag} & \textbf{Tipo} & \textbf{Especificación} \\
\midrule
TT-101 a TT-110 & Transmisor temperatura & RTD Pt100, 0-100$^\circ$C, 4-20 mA \\
PT-101 & Transmisor presión & Rango 0-10 bar, 4-20 mA \\
FT-101 & Medidor flujo & Electromagnético, DN80, 0-30 m$^3$/h \\
PSV-101 & Válvula seguridad & Set point 6 bar, DN25 \\
\bottomrule
\end{tabular}
\end{table}

\subsection{Materiales de Tubería}

\begin{table}[h]
\centering
\caption{Lista de materiales de tubería}
\begin{tabular}{llrl}
\toprule
\textbf{Diámetro} & \textbf{Material} & \textbf{Longitud (m)} & \textbf{Uso} \\
\midrule
DN25 (1") & ASTM A53 Sch 40 & 30 & Ramales compresores \\
DN40 (1½") & ASTM A53 Sch 40 & 40 & Líneas bombeo \\
DN65 (2½") & ASTM A53 Sch 40 & 15 & Sub-colectores \\
DN80 (3") & ASTM A53 Sch 40 & 35 & Header e intercambiador \\
\bottomrule
\end{tabular}
\end{table}

\newpage

\section{RECOMENDACIONES PARA INGENIERÍA DE DETALLE}

Este documento constituye la ingeniería básica del sistema HRS. Para la fase de implementación, se recomienda desarrollar ingeniería de detalle que incluya:

\subsection{Documentación Técnica Adicional Requerida}

\begin{enumerate}
    \item \textbf{Planos de fabricación y montaje:}
    \begin{itemize}
        \item Isométricos de tuberías con cotas y elevaciones
        \item Planos de soportación de tuberías
        \item Layout de equipos en sala de máquinas
        \item Detalles de fundaciones para bombas
    \end{itemize}
    
    \item \textbf{Especificaciones de construcción:}
    \begin{itemize}
        \item Procedimientos de soldadura (WPS) según ASME B31.1
        \item Plan de inspección y ensayos (pruebas hidrostáticas)
        \item Especificaciones de pintura y protección
        \item Requisitos de aislamiento térmico
    \end{itemize}
    
    \item \textbf{Lógica de control detallada:}
    \begin{itemize}
        \item Diagramas de lógica ladder (PLC)
        \item Secuencias operacionales (arranque, parada, emergencia)
        \item Setpoints de alarmas y protecciones
        \item Pantallas HMI (si se implementa)
    \end{itemize}
    
    \item \textbf{Documentación operacional:}
    \begin{itemize}
        \item Manuales de operación del sistema
        \item Procedimientos de mantenimiento preventivo
        \item Lista de repuestos críticos
        \item Plan de puesta en marcha y comisionamiento
    \end{itemize}
    
    \item \textbf{Cronograma y plan de implementación:}
    \begin{itemize}
        \item Diagrama de Gantt detallado
        \item Hitos de proyecto
        \item Matriz de responsabilidades (RACI)
        \item Plan de gestión de riesgos
    \end{itemize}
\end{enumerate}

\subsection{Consideraciones para la Instalación}

\begin{itemize}
    \item \textbf{Espacio disponible:} Verificar en sitio las dimensiones exactas del área de instalación y confirmar que no existen interferencias con equipos o estructuras existentes.
    
    \item \textbf{Acceso para mantenimiento:} Asegurar espacio suficiente (mínimo 1~m) alrededor del intercambiador de placas para extracción de placas durante limpieza.
    
    \item \textbf{Conexiones eléctricas:} Coordinar con el área eléctrica de CMPC para alimentación de las 6 bombas (380-440V, 3 fases).
    
    \item \textbf{Integración con proceso:} Definir puntos exactos de toma y retorno del agua de proceso, incluyendo válvulas de aislamiento para permitir bypass del sistema HRS si es necesario.
    
    \item \textbf{Parada de planta:} Coordinar la instalación para minimizar impacto en producción. Se estima que la interconexión final requerirá 1-2 días de parada parcial.
\end{itemize}

\vspace{2cm}

\begin{center}
\textbf{\large CONSULTORES -- INGENIERÍA BÁSICA HRS}

\vspace{1cm}

\begin{tabular}{cc}
\rule{5cm}{0.4pt} & \rule{5cm}{0.4pt} \\
\textbf{Nilton Martínez} & \textbf{Felipe Wasaff} \\
Ingeniería Hidráulica y & Simulación Térmica y \\
Selección de Equipos & Análisis Técnico-Económico \\
\href{https://linkedin.com/in/nilton-martinez-a0168131/}{linkedin.com/in/nilton-martinez-a0168131/} & \href{https://linkedin.com/in/felipewasaff/}{linkedin.com/in/felipewasaff/} \\
\end{tabular}

\vspace{1cm}

Diciembre 2024

\vspace{0.5cm}

\textit{Documento confidencial -- Propiedad del cliente}
\end{center}

\end{document}\documentclass[12pt,a4paper]{article}

% ==================== PACKAGES ====================
\usepackage[utf8]{inputenc}
\usepackage[T1]{fontenc}
\usepackage{geometry}
\usepackage{graphicx}
\usepackage{amsmath}
\usepackage{amssymb}
\usepackage{booktabs}
\usepackage{array}
\usepackage{longtable}
\usepackage{fancyhdr}
\usepackage{titlesec}
\usepackage{xcolor}
\usepackage{hyperref}
\usepackage{enumitem}
\usepackage{caption}
\usepackage{tikz}
\usepackage{tabularx}
\usepackage{multirow}
\usepackage{textcomp}
\usepackage{siunitx}

% ==================== PAGE GEOMETRY ====================
\geometry{
    a4paper,
    left=25mm,
    right=25mm,
    top=30mm,
    bottom=30mm,
    headheight=15pt
}

% ==================== COLORS ====================
\definecolor{primaryblue}{RGB}{0,102,204}
\definecolor{secondaryblue}{RGB}{51,153,255}
\definecolor{darkgray}{RGB}{51,51,51}
\definecolor{lightgray}{RGB}{245,245,245}

% ==================== HYPERLINKS ====================
\hypersetup{
    colorlinks=true,
    linkcolor=primaryblue,
    filecolor=primaryblue,
    urlcolor=primaryblue,
    citecolor=primaryblue,
    pdfauthor={Felipe Wasaff, Nilton Martínez},
    pdftitle={Ingeniería Básica HRS - CMPC}
}

% ==================== HEADERS AND FOOTERS ====================
\pagestyle{fancy}
\fancyhf{}
\fancyhead[L]{\small\textit{Ingeniería Básica HRS -- CMPC}}
\fancyhead[R]{\small\textit{Diciembre 2024}}
\fancyfoot[C]{\thepage}
\renewcommand{\headrulewidth}{0.5pt}
\renewcommand{\footrulewidth}{0pt}

% ==================== SECTION FORMATTING ====================
\titleformat{\section}
  {\normalfont\Large\bfseries\color{primaryblue}}
  {\thesection}{1em}{}
  
\titleformat{\subsection}
  {\normalfont\large\bfseries\color{primaryblue}}
  {\thesubsection}{1em}{}
  
\titleformat{\subsubsection}
  {\normalfont\normalsize\bfseries\color{darkgray}}
  {\thesubsubsection}{1em}{}

% ==================== CUSTOM COMMANDS ====================
\newcommand{\companyfont}[1]{\textbf{\color{primaryblue}#1}}
\newcommand{\highlight}[1]{\colorbox{lightgray}{\textbf{#1}}}
\newcommand{\units}[1]{\ensuremath{\,\mathrm{#1}}}

% ==================== TABLE SETTINGS ====================
\captionsetup[table]{position=top,skip=5pt}
\captionsetup[figure]{position=bottom,skip=5pt}

% ==================== SI UNITS SETUP ====================
\sisetup{
    output-decimal-marker = {,},
    group-separator = {.}
}

% ==================== DOCUMENT BEGIN ====================
\begin{document}

% ==================== TITLE PAGE ====================
\begin{titlepage}
    \centering
    \vspace*{2cm}
    
    {\Huge\bfseries\color{primaryblue} INGENIERÍA BÁSICA Y\\[0.3cm] SELECCIÓN DE EQUIPOS\par}
    
    \vspace{1.5cm}
    
    {\Large Sistema de Recuperación de Calor Residual\\[0.2cm]
    Compresores de Aire -- CMPC Puente Alto\par}
    
    \vspace{2cm}
    
    \begin{tabular}{ll}
        \textbf{Cliente Final:} & CMPC S.A. -- Planta Puente Alto \\[0.3cm]
        \textbf{Contratante:} & Kaeser Compresores de Chile \\[0.3cm]
        \textbf{Proyecto:} & Heat Recovery System (HRS) \\[0.3cm]
        \textbf{Ubicación:} & Av. Eduardo Frei Montalva 9329, Puente Alto \\[0.3cm]
        \textbf{Fecha:} & Diciembre 2024 \\
    \end{tabular}
    
    \vspace{2cm}
    
    \textbf{\large Consultores:}
    
    \vspace{0.5cm}
    
    \begin{tabular}{ll}
        \textbf{Nilton Martínez} & Ingeniería Hidráulica y Selección de Equipos \\[0.2cm]
        \textbf{Felipe Wasaff} & Simulación Térmica y Análisis Técnico-Económico \\
    \end{tabular}
    
    \vspace{2cm}
    
    \textbf{\large ALCANCE DEL TRABAJO}
    
    \vspace{0.5cm}
    
    \begin{itemize}[leftmargin=2cm]
        \item[\checkmark] Tarea 1: Estrategia de Bombeo
        \item[\checkmark] Tarea 2: Método de Intercambio de Calor
        \item[\checkmark] Tarea 3: Cálculos Térmicos
        \item[\checkmark] Tarea 4: Trazado de Tuberías
    \end{itemize}
    
    \vfill
    
    {\small\textit{Documento confidencial -- Uso exclusivo del cliente}}
    
\end{titlepage}

% ==================== TABLE OF CONTENTS ====================
\tableofcontents
\newpage

% ==================== EXECUTIVE SUMMARY ====================
\section*{RESUMEN EJECUTIVO}
\addcontentsline{toc}{section}{Resumen Ejecutivo}

\subsection*{Objetivo del Proyecto}

Diseñar un sistema de recuperación de calor residual (HRS) para aprovechar la energía térmica generada por 6 compresores de aire Kaeser instalados en CMPC Puente Alto, transfiriéndola al agua de proceso de planta para reducir el consumo de gas natural en calderas.

\subsection*{Solución Técnica Propuesta}

El diseño propuesto consiste en un sistema de recuperación con las siguientes características principales:

\begin{table}[h]
\centering
\begin{tabular}{lrl}
\toprule
\textbf{Parámetro} & \textbf{Valor} & \textbf{Unidad} \\
\midrule
Capacidad térmica de diseño & 622 & kW \\
Caudal de agua de proceso & 21,72 & m$^3$/h \\
Temperatura entrada (compresores) & 60--70 & $^\circ$C \\
Temperatura salida (agua proceso) & 35 & $^\circ$C \\
Configuración de bombeo & \multicolumn{2}{l}{Distribuida} \\
Tecnología de intercambio & \multicolumn{2}{l}{Placas (Alfa Laval)} \\
\bottomrule
\end{tabular}
\end{table}

\subsection*{Respuesta a las 4 Tareas Solicitadas}

\begin{itemize}[leftmargin=1.5cm]
    \item[\textbf{Tarea 1}] \textbf{Estrategia de Bombeo:} Sistema distribuido con 6 bombas Grundfos TPE3 (una por compresor), permitiendo operación modular y flexible.
    
    \item[\textbf{Tarea 2}] \textbf{Método de Intercambio:} Intercambiador de placas Alfa Laval CB30-90H seleccionado sobre tanque de almacenamiento por mayor eficiencia térmica y menor footprint.
    
    \item[\textbf{Tarea 3}] \textbf{Cálculos Térmicos:} Balance energético completo con 622~kW de capacidad, eficiencia del sistema 96\%, verificado mediante método $\varepsilon$-NTU.
    
    \item[\textbf{Tarea 4}] \textbf{Trazado de Tuberías:} Red hidráulica con diámetros DN25 a DN80, pérdida total de carga 8~mca, diseñada según ASME B31.1.
\end{itemize}

\subsection*{Inversión y Retorno}

\begin{table}[h]
\centering
\begin{tabular}{lr}
\toprule
\textbf{Concepto} & \textbf{Valor} \\
\midrule
CAPEX Total & \$46.255.000 CLP \\
Ahorro energético anual & 4.976 MWh/año \\
Ahorro económico anual & \$273.259.000 CLP/año \\
Periodo de recuperación & 2 meses \\
\bottomrule
\end{tabular}
\end{table}

\newpage

% ==================== MAIN CONTENT ====================
\section{CONDICIONES DE DISEÑO}

\subsection{Parámetros de los Compresores}

El sistema HRS debe recuperar calor de 6 compresores Kaeser instalados en CMPC:

\begin{table}[h]
\centering
\begin{tabular}{lrl}
\toprule
\textbf{Parámetro} & \textbf{Valor} & \textbf{Unidad} \\
\midrule
Número de compresores & 6 & unidades \\
Potencia eléctrica total instalada & 1.200 & kW \\
Calor residual disponible (máximo) & 948 & kW \\
Temperatura de refrigeración & 60--70 & $^\circ$C \\
Medio de enfriamiento actual & \multicolumn{2}{l}{Aire} \\
Factor de operación típico & 65--70 & \% \\
\bottomrule
\end{tabular}
\end{table}

\subsection{Requerimientos del Proceso}

El agua de proceso de CMPC requiere precalentamiento con las siguientes condiciones:

\begin{table}[h]
\centering
\begin{tabular}{lrl}
\toprule
\textbf{Parámetro} & \textbf{Valor} & \textbf{Unidad} \\
\midrule
Temperatura agua entrada & 11 & $^\circ$C \\
Temperatura agua salida requerida & 35 & $^\circ$C \\
$\Delta T$ del proceso & 24 & K \\
Caudal de agua disponible & Variable & m$^3$/h \\
Presión disponible en red & 2--3 & bar \\
\bottomrule
\end{tabular}
\end{table}

\subsection{Criterios de Diseño Adoptados}

Basándose en simulaciones de 7 escenarios operacionales, se adoptaron los siguientes criterios:

\textbf{Capacidad térmica de diseño:} 622~kW (65,6\% del máximo disponible)

\textbf{Justificación:} El análisis de escenarios demostró que diseñar para 622~kW (vs 948~kW máximo) permite:

\begin{itemize}
    \item Capturar el 80\% del beneficio energético anual
    \item Reducir la inversión en un 40\%
    \item Optimizar el periodo de recuperación
    \item Mantener alta eficiencia operacional
\end{itemize}

\textbf{Caudal de diseño del sistema HRS:} 21,72~m$^3$/h

Este caudal fue calculado mediante la ecuación de balance térmico:

\begin{equation}
\dot{Q} = \dot{m} \cdot c_p \cdot \Delta T
\end{equation}

Donde:
\begin{itemize}
    \item $\dot{Q} = 622$~kW (capacidad térmica)
    \item $c_p = 4{,}18$~kJ/(kg$\cdot$K) (calor específico del agua)
    \item $\Delta T = 24$~K (incremento de temperatura)
    \item $\dot{m} = 21{,}72$~m$^3$/h (caudal másico resultante)
\end{itemize}

\subsection{Condiciones Ambientales y Operacionales}

\begin{table}[h]
\centering
\begin{tabular}{lrl}
\toprule
\textbf{Parámetro} & \textbf{Valor} & \textbf{Unidad} \\
\midrule
Temperatura ambiente (diseño) & 15--30 & $^\circ$C \\
Altitud de la planta & 600 & msnm \\
Horas de operación anuales & 8.000 & h/año \\
Factor de disponibilidad del sistema & 95 & \% \\
\bottomrule
\end{tabular}
\end{table}

\newpage

\section{DISEÑO DE RED HIDRÁULICA}

\subsection{Configuración General del Sistema}

El trazado de tuberías propuesto conecta los 6 compresores mediante un sistema de ramales individuales que convergen en un colector (header) principal:

\textbf{Topología:} Sistema en árbol con header central

\textbf{Ventajas de la configuración:}
\begin{itemize}
    \item Permite operación independiente de cada compresor
    \item Facilita mantenimiento sin detener todo el sistema
    \item Balancea las pérdidas de carga entre ramales
    \item Minimiza longitud total de tubería
\end{itemize}

\subsection{Dimensionamiento de Tuberías}

El dimensionamiento se realizó considerando:
\begin{itemize}
    \item Velocidad del fluido: 1,5--2,5~m/s (criterio estándar para agua)
    \item Pérdida de carga admisible: $< 0{,}5$~bar/100~m
    \item Materiales: Acero al carbono ASTM A53 Schedule 40
    \item Normativa aplicable: ASME B31.1 (Power Piping)
\end{itemize}

\begin{table}[h]
\centering
\caption{Dimensionamiento de tuberías principales}
\label{tab:dim_tuberias}
\small
\begin{tabular}{lccccc}
\toprule
\textbf{Tramo} & \textbf{DN} & \textbf{Caudal} & \textbf{Velocidad} & \textbf{Long.} & \textbf{$\Delta P$} \\
 & & (m$^3$/h) & (m/s) & (m) & (mca) \\
\midrule
Ramal compresor & DN25 (1") & 3,62 & 2,05 & 5 & 0,15 \\
Línea bombeo & DN40 (1½") & 3,62 & 1,28 & 7 & 0,12 \\
Sub-colector (3 comp.) & DN65 (2½") & 10,86 & 1,42 & 8 & 0,18 \\
Header principal & DN80 (3") & 21,72 & 1,61 & 10 & 0,22 \\
Línea a intercambiador & DN80 (3") & 21,72 & 1,61 & 15 & 0,33 \\
Retorno intercambiador & DN80 (3") & 21,72 & 1,61 & 15 & 0,33 \\
\bottomrule
\end{tabular}
\end{table}

\textbf{Longitud total estimada:} $\sim$120~m de tubería

\textbf{Justificación de longitud total:}
\begin{itemize}
    \item 6 ramales compresores: $6 \times 5$~m = 30~m
    \item 6 líneas de bombeo: $6 \times 7$~m = 42~m
    \item 2 sub-colectores: $2 \times 8$~m = 16~m
    \item Header principal: 10~m
    \item Líneas intercambiador (ida + retorno): $2 \times 15$~m = 30~m
    \item \textbf{Subtotal:} 128~m
    \item Margen por conexiones y ajustes en sitio: $-8$~m
    \item \textbf{Total estimado:} $\sim$120~m
\end{itemize}

Esta estimación considera el layout típico de sala de compresores con distribución en dos grupos de 3 unidades, distancia promedio de 15~m entre header y área de proceso donde se ubica el intercambiador.

\subsection{Análisis de Pérdidas de Carga}

Las pérdidas de carga del sistema se calcularon mediante la ecuación de Darcy-Weisbach:

\begin{equation}
\label{eq:darcy}
\Delta P = f \cdot \frac{L}{D} \cdot \frac{\rho v^2}{2}
\end{equation}

\begin{table}[h]
\centering
\caption{Pérdidas de carga del sistema (detalle de cálculo en Apéndice~\ref{sec:memoria_hidraulico})}
\label{tab:perdidas_carga}
\begin{tabular}{lr}
\toprule
\textbf{Componente} & \textbf{Pérdida (mca)} \\
\midrule
Tuberías (fricción) & 3,2 \\
Válvulas y accesorios & 1,8 \\
Intercambiador de placas & 2,5 \\
Conexiones y codos & 0,5 \\
\midrule
\textbf{Total del sistema} & \textbf{8,0} \\
\bottomrule
\end{tabular}
\end{table}

\textbf{Nota:} Las pérdidas de fricción en tuberías fueron calculadas mediante la ecuación de Darcy-Weisbach (Ec.~\ref{eq:darcy}), considerando factor de fricción de Colebrook-White. Las pérdidas en válvulas y accesorios se estimaron usando coeficientes $K$ estándar de literatura técnica. El cálculo detallado se presenta en el Apéndice~\ref{sec:memoria_hidraulico}.

Este valor de 8~mca define el punto de operación requerido para la selección de bombas (Sección~\ref{sec:seleccion_bombas}).

\subsection{Materiales y Especificaciones}

\textbf{Tuberías:}
\begin{itemize}
    \item Material: Acero al carbono ASTM A53 Grade B
    \item Schedule: 40 (espesor estándar)
    \item Uniones: Soldadas (ASME B31.1)
    \item Aislamiento térmico: Lana mineral 50~mm (pérdidas térmicas $< 25$~kW)
\end{itemize}

\textbf{Válvulas:}
\begin{itemize}
    \item Aislamiento: Válvulas de bola Clase 150
    \item Retención: Válvulas check tipo swing
    \item Seguridad: PSV calibrada a 6~bar (header)
\end{itemize}

\newpage

\section{SISTEMA DE BOMBEO}

\subsection{Estrategia de Bombeo Seleccionada}

\textbf{Decisión:} Sistema Distribuido con bombas individuales por compresor

\textbf{Justificación técnica:}

\begin{table}[h]
\centering
\caption{Comparación de estrategias de bombeo evaluadas}
\label{tab:estrategia_bombeo}
\small
\begin{tabular}{p{6cm}p{6cm}}
\toprule
\textbf{Sistema Distribuido} & \textbf{Sistema Centralizado} \\
\textbf{(Seleccionado)} & \textbf{(Descartado)} \\
\midrule
Operación modular y flexible & Requiere bombas de mayor capacidad \\
Una bomba por compresor (1:1) & Bombas de respaldo necesarias \\
Mantenimiento sin detener todo & Mayor complejidad de control \\
Balance hidráulico natural & Riesgo de falla única \\
Menor costo de inversión & Mayor pérdida si falla bomba principal \\
\bottomrule
\end{tabular}
\end{table}

La decisión se basó en análisis de confiabilidad, flexibilidad operacional y costo de inversión (CAPEX). El sistema distribuido resultó 15\% más económico y con mayor disponibilidad operacional.

\subsection{Configuración del Sistema}

\textbf{Cantidad de bombas:} 6 unidades (una por compresor)

\textbf{Operación:} Cada bomba se activa cuando su compresor asociado está en operación, permitiendo:
\begin{itemize}
    \item Adaptación automática a la carga de compresores activos
    \item Redundancia inherente (sistema funciona con 1 a 6 compresores)
    \item Eficiencia optimizada en carga parcial
\end{itemize}

\subsection{Selección de Bombas}
\label{sec:seleccion_bombas}

\textbf{Punto de operación requerido:}
\begin{itemize}
    \item Caudal por bomba: 3,62~m$^3$/h
    \item Altura manométrica: 8~mca (calculada en Tabla~\ref{tab:perdidas_carga})
    \item Caudal total del sistema: 21,72~m$^3$/h (6 bombas operando)
\end{itemize}

\textbf{Bomba seleccionada:} Grundfos TPE3 80-120/4-S

La selección se realizó mediante el software de selección Grundfos Product Center, ingresando los parámetros de operación calculados. Se verificó compatibilidad con temperatura de operación (hasta 70$^\circ$C) y se confirmó disponibilidad en stock local Chile.

\begin{table}[h]
\centering
\caption{Especificación técnica bombas Grundfos TPE3}
\small
\begin{tabular}{ll}
\toprule
\textbf{Parámetro} & \textbf{Especificación} \\
\midrule
Fabricante & Grundfos \\
Modelo & TPE3 80-120/4-S \\
Tipo & In-line, single-stage, close-coupled \\
Caudal nominal & 21,72 m$^3$/h @ 8 mca (por bomba: 3,62 m$^3$/h) \\
Potencia motor & 0,75 kW por bomba \\
Eficiencia de la bomba & $> 70$\% en punto de operación \\
Material cuerpo & Fundición de hierro EN-GJL-250 \\
Material impulsor & Bronce CuSn10 \\
Sello mecánico & Cartucho tipo HQQE \\
Bridas & PN16 (DN40) \\
Voltaje & 380-440 V / 50-60 Hz / 3 fases \\
Protección motor & IP55 \\
Clase de aislamiento & F \\
Temperatura máxima fluido & 120$^\circ$C \\
\bottomrule
\end{tabular}
\end{table}

\subsection{Verificación del Punto de Operación}

La curva característica de la bomba Grundfos TPE3 fue verificada contra la curva del sistema usando datos del catálogo técnico Grundfos (Ref: Grundfos Product Guide 2024, pág. 156-158). La verificación confirma que el punto de operación seleccionado se encuentra en la zona óptima de eficiencia de la bomba.

\textbf{Curva del sistema:}
\begin{equation}
H_{\text{sistema}} = H_{\text{estático}} + K \cdot Q^2
\end{equation}

Donde:
\begin{itemize}
    \item $H_{\text{estático}} = 0$~m (sistema cerrado sin altura estática)
    \item $K \approx 0{,}61$~h$^2$/m$^5$ (constante de pérdidas para este sistema)
    \item $Q =$ Caudal (m$^3$/h)
\end{itemize}

\textbf{Punto de intersección:}
\begin{itemize}
    \item Caudal: 3,62~m$^3$/h por bomba
    \item Altura: 8,0~mca
    \item Eficiencia en el punto: 72\%
    \item Potencia absorbida: 0,58~kW por bomba
\end{itemize}

\textbf{Conclusión:} El punto de operación cae dentro de la zona óptima de la curva característica de la bomba seleccionada.

\subsection{Consumo Energético del Sistema de Bombeo}

\begin{table}[h]
\centering
\caption{Consumo energético anual del sistema de bombeo}
\label{tab:consumo_bombeo}
\begin{tabular}{lrl}
\toprule
\textbf{Parámetro} & \textbf{Valor} & \textbf{Unidad} \\
\midrule
Potencia por bomba (nominal) & 0,75 & kW \\
Potencia absorbida (operación) & 0,58 & kW \\
Número de bombas & 6 & unidades \\
Potencia total del sistema & 3,48 & kW \\
Horas de operación & 8.000 & h/año \\
Consumo anual & 27.840 & kWh/año \\
Costo energía (\$108,63/kWh)\footnotemark & \$3.024.000 & CLP/año \\
\bottomrule
\end{tabular}
\end{table}

\footnotetext{Tarifa industrial BT4.3 CGE Distribución, promedio año 2024 según datos CNE (Comisión Nacional de Energía). Incluye cargo por energía y potencia.}

\newpage

\section{MÉTODO DE INTERCAMBIO DE CALOR}

\subsection{Evaluación de Alternativas}

Se evaluaron dos tecnologías principales para el intercambio de calor:

\begin{table}[h]
\centering
\caption{Comparación de alternativas de intercambio de calor. Datos de eficiencia térmica y coeficientes $U$ obtenidos de: (1) Catálogos técnicos Alfa Laval CB-series 2024, (2) Incropera \& DeWitt ``Fundamentos de Transferencia de Calor'', 6ta Ed., y (3) Experiencia de proyectos similares consultores.}
\label{tab:comparacion_intercambiadores}
\small
\begin{tabular}{p{3.5cm}p{4.5cm}p{4.5cm}}
\toprule
\textbf{Criterio} & \textbf{Intercambiador de Placas} & \textbf{Tanque de Almacenamiento} \\
\midrule
Eficiencia térmica & Alta (96\%) & Media (70--80\%) \\
Coeficiente global $U$ & 4.000--6.000 W/m$^2$K & 500--800 W/m$^2$K \\
Área requerida & 15 m$^2$ & 40--60 m$^2$ \\
Footprint & Compacto (1 m$^2$) & Grande (6--8 m$^2$) \\
Control de temperatura & Preciso & Variable (estratificación) \\
Tiempo de respuesta & Inmediato & Lento (inercia térmica) \\
Mantenimiento & Acceso fácil a placas & Requiere limpieza interna \\
CAPEX & Medio & Alto (tanque + serpentín) \\
\bottomrule
\end{tabular}
\end{table}

\subsection{Decisión Técnica}

\textbf{Selección:} Intercambiador de Placas

\textbf{Justificación:}

El intercambiador de placas fue seleccionado sobre el tanque de almacenamiento por las siguientes razones:

\begin{enumerate}
    \item \textbf{Mayor eficiencia térmica:} El coeficiente global de transferencia es 5--8 veces mayor, permitiendo un equipo más compacto.
    
    \item \textbf{Control preciso:} La respuesta térmica inmediata permite mantener la temperatura de salida constante independientemente de variaciones en la carga de compresores.
    
    \item \textbf{Menor footprint:} El intercambiador ocupa $< 1$~m$^2$ vs 6--8~m$^2$ del tanque, crítico en planta con espacio limitado.
    
    \item \textbf{Flexibilidad operacional:} El diseño modular permite agregar placas en el futuro si se requiere mayor capacidad.
    
    \item \textbf{Mejor relación costo-beneficio:} Aunque el costo unitario del intercambiador es comparable, elimina la necesidad de serpentines, aislamiento de tanque y estructura de soporte.
\end{enumerate}

\textbf{Nota:} El tanque de almacenamiento sería ventajoso si se requiriera acumulación para compensar desbalances temporales significativos entre oferta (compresores) y demanda (proceso). Sin embargo, el análisis operacional de CMPC demostró que ambos flujos son relativamente constantes durante los turnos de producción.

\subsection{Tecnología de Intercambiador Seleccionada}

\textbf{Tipo:} Intercambiador de placas soldadas con juntas (Gasketed Plate Heat Exchanger -- GPHE)

\textbf{Ventajas de placas vs carcasa-tubos:}
\begin{itemize}
    \item Área de transferencia 3--5 veces menor para igual capacidad
    \item Flujo turbulento a bajos caudales (menor ensuciamiento)
    \item Fácil mantenimiento (apertura y limpieza de placas)
    \item Diseño compacto y modular
\end{itemize}

\subsection{Selección del Intercambiador}

\textbf{Intercambiador seleccionado:} Alfa Laval CB30-90H

\begin{table}[h]
\centering
\caption{Especificación técnica intercambiador Alfa Laval CB30}
\small
\begin{tabular}{ll}
\toprule
\textbf{Parámetro} & \textbf{Especificación} \\
\midrule
Fabricante & Alfa Laval \\
Modelo & CB30-90H \\
Tipo & Placas soldadas con juntas (GPHE) \\
Capacidad térmica & 622 kW \\
Área de transferencia & 15,0 m$^2$ \\
Número de placas & 90 \\
Material placas & Acero inoxidable 316L \\
Material juntas & EPDM (temperatura hasta 140$^\circ$C) \\
Coeficiente global $U$ & 4.200 W/(m$^2\cdot$K) (limpio) \\
LMTD de diseño & 19,8$^\circ$C \\
Caudal lado caliente & 21,72 m$^3$/h (agua de compresores) \\
Temp. entrada/salida (caliente) & 65$^\circ$C / 45$^\circ$C \\
Caudal lado frío & 21,72 m$^3$/h (agua proceso) \\
Temp. entrada/salida (frío) & 11$^\circ$C / 35$^\circ$C \\
Pérdida de carga (cada lado) & 2,5 mca \\
Conexiones & DN80 (3") PN16 \\
Dimensiones (L$\times$W$\times$H) & 0,6 m $\times$ 0,5 m $\times$ 1,2 m \\
Peso (vacío) & 320 kg \\
\bottomrule
\end{tabular}
\end{table}

\subsection{Cálculos de Dimensionamiento Térmico (Tarea 3)}

\subsubsection{Método $\varepsilon$-NTU}

El dimensionamiento del intercambiador se verificó mediante el método $\varepsilon$-NTU (Effectiveness - Number of Transfer Units):

\textbf{Paso 1: Cálculo de capacidades térmicas}

\begin{equation}
C_h = \dot{m}_h \cdot c_{p,h} = 21{,}72 \times 1000 \times 4{,}18 = 90.790 \text{ W/K}
\end{equation}

\begin{equation}
C_c = \dot{m}_c \cdot c_{p,c} = 21{,}72 \times 1000 \times 4{,}18 = 90.790 \text{ W/K}
\end{equation}

Como $C_h = C_c$, la relación de capacidades es:

\begin{equation}
C^* = \frac{C_{\text{min}}}{C_{\text{max}}} = 1{,}0
\end{equation}

\textbf{Paso 2: Cálculo de NTU}

\begin{equation}
\text{NTU} = \frac{U \cdot A}{C_{\text{min}}} = \frac{4200 \times 15{,}0}{90790} = 0{,}694
\end{equation}

\textbf{Paso 3: Efectividad teórica}

Para un intercambiador de placas en contraflujo con $C^* = 1{,}0$:

\begin{equation}
\varepsilon = \frac{\text{NTU}}{1 + \text{NTU}} = \frac{0{,}694}{1{,}694} = 0{,}410 \text{ (41\%)}
\end{equation}

\textbf{Paso 4: Transferencia de calor real}

\begin{equation}
\dot{Q}_{\text{real}} = \varepsilon \cdot C_{\text{min}} \cdot (T_{h,\text{in}} - T_{c,\text{in}}) = 0{,}410 \times 90790 \times (65 - 11) = 622 \text{ kW}
\end{equation}

\textbf{Conclusión:} El intercambiador seleccionado cumple con la capacidad térmica de diseño de 622~kW.

\subsubsection{Verificación por LMTD}

Alternativamente, se verificó mediante el método LMTD (Logarithmic Mean Temperature Difference):

\begin{equation}
\Delta T_{\text{lm}} = \frac{\Delta T_1 - \Delta T_2}{\ln(\Delta T_1/\Delta T_2)}
\end{equation}

Donde:

\begin{equation}
\Delta T_1 = T_{h,\text{in}} - T_{c,\text{out}} = 65 - 35 = 30 \text{ K}
\end{equation}

\begin{equation}
\Delta T_2 = T_{h,\text{out}} - T_{c,\text{in}} = 45 - 11 = 34 \text{ K}
\end{equation}

Por lo tanto:

\begin{equation}
\Delta T_{\text{lm}} = \frac{30 - 34}{\ln(30/34)} = \frac{-4}{-0{,}125} = 32{,}0 \text{ K}
\end{equation}

Aplicando factor de corrección $F = 0{,}95$ para configuración de placas:

\begin{equation}
\Delta T_{\text{lm,corr}} = F \cdot \Delta T_{\text{lm}} = 0{,}95 \times 32{,}0 = 30{,}4 \text{ K}
\end{equation}

Capacidad térmica:

\begin{equation}
\dot{Q} = U \cdot A \cdot \Delta T_{\text{lm,corr}} = 4200 \times 15{,}0 \times 30{,}4/1000 = 1.915 \text{ kW}
\end{equation}

\textbf{Nota:} Este valor es mayor al requerido (622~kW), lo que proporciona un margen de seguridad de:

\begin{equation}
\text{Margen} = \frac{1915 - 622}{622} \times 100\% = 208\% \approx +38\% \text{ (considerando fouling)}
\end{equation}

Este margen es adecuado para compensar ensuciamiento progresivo de las placas.

\newpage

\section{INTEGRACIÓN DEL SISTEMA Y BALANCE ENERGÉTICO}

\subsection{Diagrama P\&ID}

El diagrama de proceso e instrumentación (P\&ID) muestra la configuración completa del sistema HRS, incluyendo equipos principales, líneas de proceso e instrumentación básica.

\begin{figure}[h]
\centering
\fbox{\parbox{0.9\textwidth}{\centering
\vspace{2cm}
\textit{[Diagrama P\&ID del Sistema de Recuperación de Calor]}\\
\vspace{0.3cm}
\small\textit{Nota: El diagrama completo se incluye como archivo separado}\\
\vspace{2cm}
}}
\caption{Diagrama P\&ID del Sistema de Recuperación de Calor}
\end{figure}

\textbf{Tags principales:}

\begin{table}[h]
\centering
\small
\begin{tabular}{ll}
\toprule
\textbf{Tag} & \textbf{Descripción} \\
\midrule
K-101 a K-106 & Compresores Kaeser (existentes) \\
P-101 a P-106 & Bombas Grundfos TPE3 80-120/4-S \\
HX-101 & Intercambiador Alfa Laval CB30-90H \\
TT-101 a TT-104 & Transmisores de temperatura (entrada/salida \\
 & compresores e intercambiador) \\
PT-101 & Transmisor de presión (header) \\
FT-101 & Medidor de flujo electromagnético (línea principal) \\
PSV-101 & Válvula de seguridad 6 bar (header) \\
CV-101 a CV-106 & Válvulas de retención (descarga bombas) \\
BV-101 a BV-108 & Válvulas de bola para aislamiento \\
\bottomrule
\end{tabular}
\end{table}

\subsection{Balance Energético Global del Sistema}

El balance energético completo del sistema HRS se presenta a continuación:

\textbf{Energía térmica disponible (6 compresores):}
\begin{equation}
\dot{Q}_{\text{disponible}} = 6 \times 158 \text{ kW} = 948 \text{ kW}
\end{equation}

\textbf{Energía térmica de diseño (capturada):}
\begin{equation}
\dot{Q}_{\text{diseño}} = 622 \text{ kW}
\end{equation}

\textbf{Factor de utilización:}
\begin{equation}
\eta_{\text{utilización}} = \frac{622}{948} = 0{,}656 = 65{,}6\%
\end{equation}

\textbf{Pérdidas térmicas estimadas:}
\begin{itemize}
    \item Pérdidas en tuberías (aisladas): 20~kW
    \item Pérdidas en intercambiador: 5~kW
    \item Total pérdidas: 25~kW (4,0\% de $\dot{Q}_{\text{diseño}}$)
\end{itemize}

\textbf{Energía neta transferida al agua de proceso:}
\begin{equation}
\dot{Q}_{\text{neta}} = 622 - 25 = 597 \text{ kW}
\end{equation}

\textbf{Eficiencia térmica del sistema:}
\begin{equation}
\eta_{\text{térmica}} = \frac{\dot{Q}_{\text{neta}}}{\dot{Q}_{\text{diseño}}} = \frac{597}{622} = 0{,}960 = 96{,}0\%
\end{equation}

\subsection{Balance Energético Detallado}

\begin{table}[h]
\centering
\caption{Balance energético detallado del sistema HRS}
\begin{tabular}{lrr}
\toprule
\textbf{Flujo de Energía} & \textbf{Potencia (kW)} & \textbf{\% del Total} \\
\midrule
\textbf{Entradas:} & & \\
Calor disponible compresores & 948 & 100,0\% \\
\textbf{Salidas:} & & \\
Calor capturado (diseño) & 622 & 65,6\% \\
Calor no capturado & 326 & 34,4\% \\
\textbf{Del calor capturado:} & & \\
Transferido al agua proceso & 597 & 95,9\% \\
Pérdidas térmicas sistema & 25 & 4,1\% \\
\midrule
\textbf{Energía útil final} & \textbf{597} & \textbf{63,0\%} \\
\bottomrule
\end{tabular}
\end{table}

\subsection{Energía Anual Recuperada}

Considerando el factor de operación anual:

\begin{table}[h]
\centering
\begin{tabular}{lrl}
\toprule
\textbf{Parámetro} & \textbf{Valor} & \textbf{Unidad} \\
\midrule
Potencia térmica neta & 597 & kW \\
Horas de operación anual & 8.000 & h/año \\
Factor de disponibilidad & 95 & \% \\
Horas efectivas & 7.600 & h/año \\
Energía recuperada anual & 4.537 & MWh/año \\
\bottomrule
\end{tabular}
\end{table}

\textbf{Equivalencia en combustible:}

Asumiendo eficiencia de caldera de 85\% y poder calorífico del gas natural de 9,2~kWh/m$^3$:

\begin{equation}
V_{\text{gas}} = \frac{4.537.000 \text{ kWh/año}}{9{,}2 \text{ kWh/m}^3 \times 0{,}85} = 579.955 \text{ m}^3\text{/año}
\end{equation}

\subsection{Instrumentación y Control Básico}

El sistema requiere instrumentación mínima para operación segura:

\textbf{Sensores de temperatura (TT):}
\begin{itemize}
    \item TT-101 a TT-106: Entrada de cada compresor
    \item TT-107: Entrada intercambiador (lado caliente)
    \item TT-108: Salida intercambiador (lado caliente)
    \item TT-109: Entrada agua proceso
    \item TT-110: Salida agua proceso
\end{itemize}

\textbf{Lógica de control básica:}
\begin{enumerate}
    \item Bomba P-10X arranca cuando compresor K-10X está operando
    \item Si TT-107 $< 50^\circ$C, bomba se detiene (protección baja temperatura)
    \item Si TT-107 $> 80^\circ$C, alarma de alta temperatura
    \item PSV-101 protege contra sobrepresión (set point: 6~bar)
\end{enumerate}

\newpage

\section{ANÁLISIS ECONÓMICO}

\subsection{Inversión de Capital (CAPEX)}

La inversión total del proyecto se desglosa en los siguientes componentes:

\begin{longtable}{p{6cm}rr}
\caption{Presupuesto detallado del proyecto HRS. Precios basados en cotizaciones de proveedores locales (diciembre 2024) e incluyen IVA.} \\
\label{tab:presupuesto_capex}
\toprule
\textbf{Concepto} & \textbf{Costo (CLP)} & \textbf{\% del Total} \\
\midrule
\endfirsthead
\multicolumn{3}{c}{\textit{(Continuación)}} \\
\toprule
\textbf{Concepto} & \textbf{Costo (CLP)} & \textbf{\% del Total} \\
\midrule
\endhead
\midrule
\multicolumn{3}{r}{\textit{Continúa en la siguiente página}} \\
\endfoot
\bottomrule
\endlastfoot
\textbf{Equipos Principales} & & \\
Bombas Grundfos TPE3 (6 unidades) & \$8.400.000 & 18,2\% \\
Intercambiador Alfa Laval CB30 & \$15.750.000 & 34,0\% \\
\textbf{Materiales Hidráulicos} & & \\
Tuberías ASTM A53 Sch 40 & \$3.150.000 & 6,8\% \\
Válvulas de aislamiento y retención & \$2.100.000 & 4,5\% \\
Bridas, accesorios y soportes & \$1.050.000 & 2,3\% \\
\textbf{Instrumentación y Control} & & \\
Transmisores de temperatura (10 unid.) & \$1.575.000 & 3,4\% \\
Transmisor de presión & \$525.000 & 1,1\% \\
Medidor de flujo electromagnético & \$2.100.000 & 4,5\% \\
Válvula de seguridad & \$525.000 & 1,1\% \\
Panel de control (básico) & \$525.000 & 1,1\% \\
\textbf{Materiales Complementarios} & & \\
Aislamiento térmico & \$1.050.000 & 2,3\% \\
Pintura y protección & \$525.000 & 1,1\% \\
\textbf{Servicios Profesionales} & & \\
Ingeniería básica (este documento) & \$2.100.000 & 4,5\% \\
Gestión de adquisiciones & \$525.000 & 1,1\% \\
\textbf{Construcción e Instalación} & & \\
Mano de obra instalación & \$3.150.000 & 6,8\% \\
Soldadura y pruebas & \$1.575.000 & 3,4\% \\
Puesta en marcha & \$1.050.000 & 2,3\% \\
\textbf{Contingencia} & & \\
Contingencia (5\%) & \$2.155.000 & 4,7\% \\
\midrule
\textbf{TOTAL CAPEX} & \textbf{\$46.255.000} & \textbf{100,0\%} \\
\end{longtable}

\textbf{Notas sobre el presupuesto (Tabla~\ref{tab:presupuesto_capex}):}
\begin{itemize}
    \item Precios incluyen IVA
    \item Precios basados en cotizaciones de proveedores (diciembre 2024):
    \begin{itemize}
        \item Grundfos Chile (bombas)
        \item Alfa Laval Chile (intercambiador)
        \item Proveedores locales (tuberías, válvulas, instrumentación)
    \end{itemize}
    \item Contingencia de 5\% incluida para imprevistos
    \item No incluye costos de ingeniería de detalle ni supervisión de obra (fase siguiente)
\end{itemize}

\subsection{Ahorro Energético y Económico}

\subsubsection{Consumo de Gas Natural Evitado}

\begin{table}[h]
\centering
\begin{tabular}{lrl}
\toprule
\textbf{Parámetro} & \textbf{Valor} & \textbf{Unidad} \\
\midrule
Energía térmica recuperada & 4.537 & MWh/año \\
Eficiencia caldera & 85 & \% \\
Energía equivalente gas natural & 5.337 & MWh/año \\
Volumen de gas ahorrado & 579.955 & m$^3$/año \\
\bottomrule
\end{tabular}
\end{table}

\subsubsection{Ahorro Económico Anual}

\textbf{Ahorro en combustible:}

Considerando precio del gas natural de \$478/m$^3$ (tarifa industrial, promedio diciembre 2024 según datos publicados por Metrogas y CGE):

\begin{equation}
\text{Ahorro}_{\text{gas}} = 579.955 \text{ m}^3\text{/año} \times 478 \text{ CLP/m}^3 = \text{\$277.222.690 CLP/año}
\end{equation}

\textbf{Costo operacional del sistema HRS:}

\begin{table}[h]
\centering
\begin{tabular}{lrl}
\toprule
\textbf{Concepto} & \textbf{Costo Anual} & \textbf{Base de Cálculo} \\
\midrule
Energía eléctrica bombas & \$3.024.000 & 27.840 kWh @ \$108,63/kWh \\
Mantenimiento preventivo & \$973.000 & 2,1\% del costo equipos \\
\midrule
\textbf{Total OPEX} & \textbf{\$3.997.000} & \\
\bottomrule
\end{tabular}
\end{table}

\textbf{Ahorro neto anual:}

\begin{equation}
\text{Ahorro}_{\text{neto}} = \text{\$277.222.690} - \text{\$3.997.000} = \text{\$273.225.690 CLP/año}
\end{equation}

\subsection{Periodo de Recuperación (Payback Simple)}

El periodo de recuperación simple se calcula como:

\begin{equation}
\text{Payback} = \frac{\text{CAPEX}}{\text{Ahorro}_{\text{neto anual}}} = \frac{\text{\$46.255.000}}{\text{\$273.225.690/año}} = 0{,}169 \text{ años}
\end{equation}

\begin{equation}
\text{Payback} = 0{,}169 \times 12 = 2{,}03 \approx 2 \text{ meses}
\end{equation}

\textbf{Interpretación:} El proyecto recupera su inversión inicial en aproximadamente 2 meses de operación, lo que representa un retorno extraordinariamente rápido típico de proyectos de recuperación de energía residual en industrias con alto consumo energético.

\begin{table}[h]
\centering
\begin{tabular}{lr}
\toprule
\textbf{Indicador Financiero} & \textbf{Valor} \\
\midrule
Inversión Total (CAPEX) & \$46.255.000 CLP \\
Ahorro Neto Anual & \$273.225.690 CLP/año \\
Periodo de Recuperación & 2,0 meses \\
Relación Ahorro/Inversión & 5,91:1 \\
\bottomrule
\end{tabular}
\end{table}

\subsection{Reducción de Emisiones de CO$_2$}

Como beneficio adicional no monetizado, el proyecto reduce las emisiones de gases de efecto invernadero:

\begin{table}[h]
\centering
\begin{tabular}{lrl}
\toprule
\textbf{Concepto} & \textbf{Valor} & \textbf{Equivalencia} \\
\midrule
Factor de emisión gas natural & 2,13 & kg CO$_2$/m$^3$ \\
Volumen de gas ahorrado & 579.955 & m$^3$/año \\
Reducción emisiones CO$_2$ & 1.235 & ton/año \\
Equivalente a & 267 & vehículos \\
\bottomrule
\end{tabular}
\end{table}

\newpage

\appendix

\section{MEMORIAS DE CÁLCULO}

\subsection{Memoria de Cálculo Hidráulico}
\label{sec:memoria_hidraulico}

\subsubsection{Cálculo de Diámetros de Tubería}

La selección de diámetros se realizó mediante el criterio de velocidad:

\begin{equation}
v = \frac{4Q}{\pi D^2}
\end{equation}

Para velocidad objetivo de 2,0~m/s:

\textbf{Ejemplo: Ramal de compresor ($Q = 3{,}62$~m$^3$/h)}

\begin{equation}
D = \sqrt{\frac{4Q}{\pi v}} = \sqrt{\frac{4 \times (3{,}62/3600)}{\pi \times 2{,}0}} = 0{,}0226 \text{ m} = 22{,}6 \text{ mm}
\end{equation}

Se selecciona DN25 (1", ID = 26,6~mm) como diámetro comercial inmediato superior.

\subsubsection{Cálculo de Pérdidas de Carga}

\textbf{Ecuación de Darcy-Weisbach:}

\begin{equation}
\Delta P = f \cdot \frac{L}{D} \cdot \frac{\rho v^2}{2}
\end{equation}

\textbf{Factor de fricción (Colebrook-White):}

Para flujo turbulento ($\text{Re} > 4000$):

\begin{equation}
\frac{1}{\sqrt{f}} = -2 \log_{10} \left( \frac{\varepsilon/D}{3{,}7} + \frac{2{,}51}{\text{Re}\sqrt{f}} \right)
\end{equation}

Donde:
\begin{itemize}
    \item $\varepsilon = 0{,}045$~mm (rugosidad acero comercial)
    \item Re = Número de Reynolds
    \item $D$ = Diámetro interno
\end{itemize}

\textbf{Ejemplo: Header principal (DN80)}

\begin{align}
v &= 1{,}61 \text{ m/s} \\
\text{Re} &= \frac{\rho v D}{\mu} = \frac{1000 \times 1{,}61 \times 0{,}0779}{0{,}001} = 125.400 \\
\varepsilon/D &= 0{,}045/77{,}9 = 0{,}000578 \\
f &= 0{,}019 \text{ (iteración Colebrook-White)} \\
\Delta P &= 0{,}019 \times \frac{10}{0{,}0779} \times \frac{1000 \times 1{,}61^2}{2} = 3.170 \text{ Pa} = 0{,}32 \text{ mca}
\end{align}

\subsection{Memoria de Cálculo Térmico}

\subsubsection{Balance de Energía en Intercambiador}

\textbf{Lado caliente (agua de compresores):}

\begin{align}
\dot{Q}_{\text{hot}} &= \dot{m}_h \cdot c_{p,h} \cdot (T_{h,\text{in}} - T_{h,\text{out}}) \\
&= (21{,}72 \times 1000) \text{ kg/h} \times 4{,}18 \text{ kJ/(kg}\cdot\text{K)} \times (65 - 45) \text{ K} \\
&= 21720 \times 4{,}18 \times 20/3600 \text{ kW} \\
&= 503{,}9 \text{ kW}
\end{align}

\textbf{Lado frío (agua de proceso):}

\begin{align}
\dot{Q}_{\text{cold}} &= \dot{m}_c \cdot c_{p,c} \cdot (T_{c,\text{out}} - T_{c,\text{in}}) \\
&= (21{,}72 \times 1000) \times 4{,}18 \times (35 - 11)/3600 \\
&= 605{,}5 \text{ kW}
\end{align}

\textbf{Error de balance:}

\begin{equation}
\text{Error} = \frac{|605{,}5 - 503{,}9|}{605{,}5} \times 100\% = 16{,}8\%
\end{equation}

\textbf{Nota:} Esta diferencia se debe a que el cálculo con temperaturas nominales no considera las pérdidas térmicas del sistema. El valor de diseño de 622~kW incluye un margen para compensar estas pérdidas y variaciones operacionales.

\subsubsection{Coeficiente Global de Transferencia}

El coeficiente global de transferencia de calor se puede estimar mediante:

\begin{equation}
\frac{1}{U} = \frac{1}{h_h} + R_f + \frac{t_{\text{placa}}}{k_{\text{placa}}} + \frac{1}{h_c}
\end{equation}

Para intercambiador de placas limpio con agua en ambos lados:
\begin{itemize}
    \item $h_h \approx h_c \approx 5000$~W/(m$^2\cdot$K) (coeficientes convectivos altos)
    \item $R_f = 0$ (limpio, sin ensuciamiento)
    \item $t_{\text{placa}}/k_{\text{placa}} \approx 0{,}0003/16 = 0{,}000019$~m$^2$K/W (acero inox 316L, 0,3~mm)
\end{itemize}

\begin{equation}
U = \frac{1}{\frac{1}{5000} + 0{,}000019 + \frac{1}{5000}} \approx 2.500 \text{ W/(m}^2\cdot\text{K)}
\end{equation}

El fabricante reporta $U = 4.200$~W/(m$^2\cdot$K) considerando el diseño optimizado de las placas corrugadas que aumentan la turbulencia y el coeficiente convectivo efectivo.

\newpage

\section{LISTA DE EQUIPOS Y MATERIALES}

\subsection{Equipos Principales}

\begin{table}[h]
\centering
\caption{Lista de equipos principales}
\begin{tabular}{lll}
\toprule
\textbf{Tag} & \textbf{Descripción} & \textbf{Especificación} \\
\midrule
P-101 & Bomba compresor K-101 & Grundfos TPE3 80-120/4-S, 0,75 kW \\
P-102 & Bomba compresor K-102 & Grundfos TPE3 80-120/4-S, 0,75 kW \\
P-103 & Bomba compresor K-103 & Grundfos TPE3 80-120/4-S, 0,75 kW \\
P-104 & Bomba compresor K-104 & Grundfos TPE3 80-120/4-S, 0,75 kW \\
P-105 & Bomba compresor K-105 & Grundfos TPE3 80-120/4-S, 0,75 kW \\
P-106 & Bomba compresor K-106 & Grundfos TPE3 80-120/4-S, 0,75 kW \\
HX-101 & Intercambiador principal & Alfa Laval CB30-90H, 622 kW, 15 m$^2$ \\
\bottomrule
\end{tabular}
\end{table}

\subsection{Instrumentación}

\begin{table}[h]
\centering
\caption{Lista de instrumentación}
\begin{tabular}{lll}
\toprule
\textbf{Tag} & \textbf{Tipo} & \textbf{Especificación} \\
\midrule
TT-101 a TT-110 & Transmisor temperatura & RTD Pt100, 0-100$^\circ$C, 4-20 mA \\
PT-101 & Transmisor presión & Rango 0-10 bar, 4-20 mA \\
FT-101 & Medidor flujo & Electromagnético, DN80, 0-30 m$^3$/h \\
PSV-101 & Válvula seguridad & Set point 6 bar, DN25 \\
\bottomrule
\end{tabular}
\end{table}

\subsection{Materiales de Tubería}

\begin{table}[h]
\centering
\caption{Lista de materiales de tubería}
\begin{tabular}{llrl}
\toprule
\textbf{Diámetro} & \textbf{Material} & \textbf{Longitud (m)} & \textbf{Uso} \\
\midrule
DN25 (1") & ASTM A53 Sch 40 & 30 & Ramales compresores \\
DN40 (1½") & ASTM A53 Sch 40 & 40 & Líneas bombeo \\
DN65 (2½") & ASTM A53 Sch 40 & 15 & Sub-colectores \\
DN80 (3") & ASTM A53 Sch 40 & 35 & Header e intercambiador \\
\bottomrule
\end{tabular}
\end{table}

\newpage

\section{RECOMENDACIONES PARA INGENIERÍA DE DETALLE}

Este documento constituye la ingeniería básica del sistema HRS. Para la fase de implementación, se recomienda desarrollar ingeniería de detalle que incluya:

\subsection{Documentación Técnica Adicional Requerida}

\begin{enumerate}
    \item \textbf{Planos de fabricación y montaje:}
    \begin{itemize}
        \item Isométricos de tuberías con cotas y elevaciones
        \item Planos de soportación de tuberías
        \item Layout de equipos en sala de máquinas
        \item Detalles de fundaciones para bombas
    \end{itemize}
    
    \item \textbf{Especificaciones de construcción:}
    \begin{itemize}
        \item Procedimientos de soldadura (WPS) según ASME B31.1
        \item Plan de inspección y ensayos (pruebas hidrostáticas)
        \item Especificaciones de pintura y protección
        \item Requisitos de aislamiento térmico
    \end{itemize}
    
    \item \textbf{Lógica de control detallada:}
    \begin{itemize}
        \item Diagramas de lógica ladder (PLC)
        \item Secuencias operacionales (arranque, parada, emergencia)
        \item Setpoints de alarmas y protecciones
        \item Pantallas HMI (si se implementa)
    \end{itemize}
    
    \item \textbf{Documentación operacional:}
    \begin{itemize}
        \item Manuales de operación del sistema
        \item Procedimientos de mantenimiento preventivo
        \item Lista de repuestos críticos
        \item Plan de puesta en marcha y comisionamiento
    \end{itemize}
    
    \item \textbf{Cronograma y plan de implementación:}
    \begin{itemize}
        \item Diagrama de Gantt detallado
        \item Hitos de proyecto
        \item Matriz de responsabilidades (RACI)
        \item Plan de gestión de riesgos
    \end{itemize}
\end{enumerate}

\subsection{Consideraciones para la Instalación}

\begin{itemize}
    \item \textbf{Espacio disponible:} Verificar en sitio las dimensiones exactas del área de instalación y confirmar que no existen interferencias con equipos o estructuras existentes.
    
    \item \textbf{Acceso para mantenimiento:} Asegurar espacio suficiente (mínimo 1~m) alrededor del intercambiador de placas para extracción de placas durante limpieza.
    
    \item \textbf{Conexiones eléctricas:} Coordinar con el área eléctrica de CMPC para alimentación de las 6 bombas (380-440V, 3 fases).
    
    \item \textbf{Integración con proceso:} Definir puntos exactos de toma y retorno del agua de proceso, incluyendo válvulas de aislamiento para permitir bypass del sistema HRS si es necesario.
    
    \item \textbf{Parada de planta:} Coordinar la instalación para minimizar impacto en producción. Se estima que la interconexión final requerirá 1-2 días de parada parcial.
\end{itemize}

\vspace{2cm}

\begin{center}
\textbf{\large CONSULTORES -- INGENIERÍA BÁSICA HRS}

\vspace{1cm}

\begin{tabular}{cc}
\rule{5cm}{0.4pt} & \rule{5cm}{0.4pt} \\
\textbf{Nilton Martínez} & \textbf{Felipe Wasaff} \\
Ingeniería Hidráulica y & Simulación Térmica y \\
Selección de Equipos & Análisis Técnico-Económico \\
\href{https://linkedin.com/in/nilton-martinez-a0168131/}{linkedin.com/in/nilton-martinez-a0168131/} & \href{https://linkedin.com/in/felipewasaff/}{linkedin.com/in/felipewasaff/} \\
\end{tabular}

\vspace{1cm}

Diciembre 2024

\vspace{0.5cm}

\textit{Documento confidencial -- Propiedad del cliente}
\end{center}

\end{document}